\section{The primary arithmetic object and the topology actually specified}
\label{sec:cq-source}
The starting point is Connes--Consani--Moscovici,
\emph{Zeta zeros and prolate wave operators: Semilocal adelic operators},
Section 3.6, proof-PDF pages 17--19, especially Lemma 3.3,
equations (3.16)--(3.22), and Proposition 3.6
\cite{cqCCM}. The preceding arithmetic map is
\[
 \mathcal E f(x)=x^{1/2}\sum_{n\geq1}f(nx),\qquad
 \mathcal F_\mu g(s)=\int_0^\infty g(x)x^{-is}\,\frac{dx}{x}.
\]
The input consists of even Schwartz functions with
$f(0)=\widehat f(0)=0$. The paper's two Hermite families are
\[
 \psi_\ell^+=h_{4\ell}-\frac{h_{4\ell}(0)}{h_0(0)}h_0,
 \qquad
 \psi_\ell^-=-h_{4\ell+2}
       +\frac{h_{4\ell+2}(0)}{h_2(0)}h_2\quad(\ell\geq1).
\]
These signs and the minus sign in the Mellin exponent are retained.
The completed function is
\begin{equation}\label{eq:cq-Xi}
 \Xi(s)=\xi(\tfrac12+is),\qquad
 \xi(z)=\tfrac12z(z-1)\pi^{-z/2}\Gamma(z/2)\zeta(z),
 \qquad \Xi(s)=8H_0(2s).
\end{equation}
The paper proves
\begin{equation}\label{eq:cq-image}
 \mathcal F_\mu\mathcal E\psi_\ell^\pm=R_\ell^\pm\Xi,
 \qquad P_\ell^\pm=-\tfrac12(\tfrac14+s^2)R_\ell^\pm.
\end{equation}
It also states immediately before Proposition 3.6 that these give all
polynomial multiples of $\Xi$. Here is the precise algebra behind the
last assertion. The top summands of (3.19) and (3.20) are nonzero;
their degrees are $2\ell$ and $2\ell+1$. Division by
$-\tfrac12(\tfrac14+s^2)$ gives degrees $2\ell-2$ and
$2\ell-1$, respectively. The even and odd parities remain.
Thus the ordered family $R_1^+,R_1^-,R_2^+,R_2^-,\ldots$
has one polynomial of each degree, with a nonzero leading coefficient.
Subtracting the appropriate multiple of the polynomial of largest
degree, and iterating down to degree zero, expresses every polynomial
as a finite linear combination. Linear independence follows by the
largest-degree coefficient of a finite relation. Consequently
\begin{equation}\label{eq:cq-N}
 \cqN=\overline{\operatorname{span}
 \{\mathcal F_\mu\mathcal E\psi_\ell^\pm\}}^{\,\tau}
 =\overline{\cqC[s]\Xi}^{\,\tau},\qquad
 \cqQ=(\cqE,\tau)/\cqN.
\end{equation}
Here $\tau$ denotes precisely the source's topology. No equality of
this closed space with a principal algebraic ideal is being asserted.
Multiplication $M f(s)=sf(s)$ is continuous in the topological ring;
it preserves the dense generating subspace, and hence $\cqN$.
Its induced operator $S$ is the operator of Proposition 3.6.
We use the source's spectral assertion as an explicitly imported theorem:
\begin{equation}\label{eq:cq-source-spectrum}
 \cqSpec(S)=\{\lambda\in\cqC:
             \zeta(\tfrac12+i\lambda)=0
             \text{ at a nontrivial zero}\}.
\end{equation}
Spectrum here means failure of a continuous two-sided inverse.
No proof of the entire source spectral theorem is claimed as a new
result of this manuscript.

There is a material source limitation. Page 19 calls $\cqE$ the
Hadamard topological ring of entire functions of order at most one;
it does not state seminorms or a convergence criterion. The cited
Connes--Marcolli Proposition 2.24, printed pages 379--381,
proves the polynomial-image assertion and explicitly discusses
different function-space realizations \cite{cqCM}. It does not supply
that missing specification. We therefore retain $\tau$ in every
source quotient. Below, all assertions depending on another topology
name and define that topology. The comparison with the source is an
exact quotient diagram, rather than an identification by terminology.
The bounded edition audit gives the same limitation in the primary
arXiv v1 source, lines 683--704 (Proposition 3.9), and in v2,
\texttt{mainc2m24fine.tex}, lines 668--689 (Proposition 3.6):
each states the order class, polynomial trace image, and spectral
set, without defining its topology or algebraic multiplicities.
These exact edition locators supplement the author-proof locator
in \cite{cqCCM}; no claim of byte identity with a publisher's final
subscription text is required.

\section{A globally defined heat map on the original underlying functions}
\label{sec:cq-heatspace}
For an entire function $f$, let
$M_f(r)=\max_{|s|\leq r}|f(s)|$ and define its order by
\[
 \limsup_{r\to\infty}
       \frac{\log\log\max(e,M_f(r))}{\log r}.
\]
The zero function is included. For $1<p<2$, $a>0$ set
\[
 \|f\|_{p,a}=\sup_{s\in\cqC}|f(s)|e^{-a|s|^p}.
\]
The underlying set $\cqE$ is exactly the set where all these seminorms
are finite. Indeed order at most one gives, for $1<q<p$,
$M_f(r)\leq\exp(r^q)$ for sufficiently large $r$; the expression
$r^q-ar^p$ is bounded above. Conversely finiteness gives order at
most $p$ for every $p>1$.
Write $\gamma$ for this explicitly defined growth topology.
It is enough to take $p=1+1/j$, $j\geq2$, and $a=1/k$, $k\geq1$:
given $p>1,a>0$, select one of these $p_j<p$ and use the boundedness
of $|s|^{p_j}/k-a|s|^p$. This proves equivalence of the countable
family with all displayed seminorms. A sequence Cauchy in these
seminorms converges uniformly on compact sets to an entire function.
Taking pointwise limits in each weighted Cauchy inequality gives
convergence in that seminorm and finiteness of its norm. Thus
$(\cqE,\gamma)$ is a complete metrizable locally convex space.

\begin{cqlemma}[Derivative bounds and the heat group]\label{thm:cq-heat}
Put $D=\cqD$. Differentiation and multiplication by $s$ map the
original underlying set $\cqE$ into itself. For $m\geq1$,
\begin{equation}\label{eq:cq-cauchy}
 \|D^mf\|_{p,a}\leq
 m!\left(\frac{epa}{m}\right)^{m/p}
 \|f\|_{p,a/2^{p-1}}.
\end{equation}
For every $t\in\cqC$, the series
\begin{equation}\label{eq:cq-U}
 U_tf=\sum_{k=0}^\infty\frac{(-t/4)^k}{k!}D^{2k}f
\end{equation}
converges in $\gamma$, locally uniformly in $t$; $U_t$ is a
continuous automorphism of $(\cqE,\gamma)$ and
$U_tU_u=U_{t+u}$, $U_t^{-1}=U_{-t}$.
\end{cqlemma}
\begin{proof}
Cauchy's circle formula on a circle of radius $R$ about $s$ gives
\[
 |D^mf(s)|\leq m!R^{-m}\max_{|z-s|=R}|f(z)|.
\]
For $b=a/2^{p-1}$ the inequality
$|s+z|^p\leq2^{p-1}(|s|^p+|z|^p)$ bounds the last expression,
after multiplication by $e^{-a|s|^p}$, by
$m!R^{-m}e^{aR^p}\|f\|_{p,b}$. Choosing
$R=(m/(ap))^{1/p}$ proves \eqref{eq:cq-cauchy}.
For multiplication use
\[
 \|Mf\|_{p,a}\leq\sup_{r\geq0}re^{-ar^p/2}\|f\|_{p,a/2}.
\]
The $k$th heat summand is at most
\[
 \frac{(|t|/4)^k}{k!}(2k)!
       \left(\frac{epa}{2k}\right)^{2k/p}\|f\|_{p,b}.
\]
Since $(2k)!/k!\leq(2k)^k$, its scalar factor is at most
$[C_{p,a}|t|k^{1-2/p}]^k$, whose $k$th root tends to zero.
This proves summability, continuity, and locally uniform convergence
of every time-differentiated series by the same estimate with the
additional polynomial factors in $k$. Commutation with $D$ follows
from continuity of $D$. In the double composition series group
terms with $n=k+\ell$. The sum of the absolute scalar coefficients
at $D^{2n}f$ is $( (|t|+|u|)/4)^n/n!$; the preceding estimate
justifies regrouping. The binomial theorem gives $U_{t+u}$.
The value at zero is the identity, proving the inverse.
\end{proof}

\begin{cqlemma}[Polynomial approximation]\label{thm:cq-density}
The Taylor polynomials of any $f\in\cqE$ converge to $f$ in $\gamma$.
Multiplication of two elements is continuous in $\gamma$.
\end{cqlemma}
\begin{proof}
If $f=\sum c_ns^n$, Cauchy's coefficient bound with parameter
$b<a$ gives, for $n\geq1$,
\[
 |c_n|\leq\|f\|_{p,b}(epb/n)^{n/p}.
\]
Also $\|s^n\|_{p,a}=(n/(epa))^{n/p}$. Therefore
$\|c_ns^n\|_{p,a}\leq\|f\|_{p,b}(b/a)^{n/p}$.
The tail is bounded by a convergent geometric series.
Finally $\|fg\|_{p,a}\leq\|f\|_{p,a/2}\|g\|_{p,a/2}$.
\end{proof}

Retain the actual real heat family, rather than substitute a polynomial:
\begin{align}\label{eq:cq-realheat}
 \Phi(u)&=\sum_{n\geq1}
 (2\pi^2n^4e^{9u}-3\pi n^2e^{5u})e^{-\pi n^2e^{4u}},\notag\\
 Z_t(s)&=8H_t(2s)=8\int_0^\infty e^{tu^2}\Phi(u)\cos(2su)\,du,
 &\partial_tZ_t=-\tfrac14D^2Z_t,
 &Z_0=\Xi.
\end{align}
These are the Rodgers--Tao conventions, Section 1 \cite{cqRT}.
For clarity the analytic exchanges in this transport can be proved
directly. Since $n^2e^{4u}\geq(n^2+e^{4u})/2$, the absolutely
convergent sums $\sum n^je^{-\pi n^2/2}$ give constants $C,c>0$
such that $|\Phi(u)|\leq Ce^{9u}e^{-ce^{4u}}$ for $u\geq0$.
For $|t|\leq T$ and $|s|\leq R$, every derivative of the integrand
is bounded by a polynomial in $u$ times
$C\exp(Tu^2+(2R+9)u-ce^{4u})$, which is integrable.
The exponential series in $t$ can therefore be integrated absolutely.
Because $D^{2k}\cos(2su)=(-1)^k(2u)^{2k}\cos(2su)$, it gives
exactly $Z_t=U_t\Xi$ with the coefficient $(-t/4)^k/k!$.
In particular $Z_t$ lies in the original order-at-most-one set for
every complex $t$, by Lemma~\ref{thm:cq-heat}; it is never identically
zero because $U_t$ is invertible and $\Xi(0)\ne0$.

\section{Transport of the actual closed arithmetic image}
\label{sec:cq-moving}
Define $\tau_t$ on the same underlying vector space $\cqE$ by
declaring $U_t:(\cqE,\tau)\longrightarrow(\cqE,\tau_t)$ a
homeomorphism. This definition specifies every open set:
$V$ is $\tau_t$-open exactly when $U_{-t}V$ is $\tau$-open.
It uses the globally defined bijection just proved and does not require
unproved continuity in $\tau$. Set
\begin{equation}\label{eq:cq-Qt}
 N_t=U_t\cqN
 =\overline{\operatorname{span}\{U_t(R_\ell^\pm\Xi)\}}^{\,\tau_t},
 \qquad Q_t=(\cqE,\tau_t)/N_t.
\end{equation}
The closure equality follows because homeomorphisms carry closure to
closure. In particular $N_t$ is closed. The quotient map
\begin{equation}\label{eq:cq-Phi}
 \mathcal U_t:\cqQ\longrightarrow Q_t,\qquad [f]\longmapsto[U_tf]
\end{equation}
is well defined, bijective, and continuous, with continuous inverse
$[g]\mapsto[U_{-t}g]$. Indeed changing $f$ by $n\in\cqN$
changes $U_tf$ by $U_tn\in N_t$, and the inverse gives necessity;
the quotient universal property proves both continuity statements.

\begin{cqtheorem}[Exact arithmetic heat intertwiner]\label{thm:cq-moving}
On the entire-function set,
\begin{equation}\label{eq:cq-B}
 B_t=U_tMU_{-t}=M-\tfrac t2D.
\end{equation}
This operator preserves $N_t$, is continuous in $\tau_t$, and induces
$S_t$ on $Q_t$. It satisfies
$S_t\mathcal U_t=\mathcal U_tS$ and
\begin{equation}\label{eq:cq-specmove}
 \cqSpec(S_t)=\cqSpec(S)
 =\{\lambda:\zeta(\tfrac12+i\lambda)=0\text{ nontrivially}\}.
\end{equation}
The time-dependent trace functions are exactly
$U_t(R_\ell^\pm\Xi)=R_\ell^\pm(B_t)Z_t$.
\end{cqtheorem}
\begin{proof}
The product rule gives $D^{2k}Mf=MD^{2k}f+2kD^{2k-1}f$.
Insert this identity in the convergent heat series to obtain
$U_tM=(M-(t/2)D)U_t$, proving \eqref{eq:cq-B}.
If $n\in\cqN$, then $B_tU_tn=U_tMn\in N_t$.
Continuity is conjugation of the continuous $M$ by the defining
homeomorphism. The induced identity follows on representatives.
If $T$ is a continuous inverse of $S-\lambda$, then
$\mathcal U_tT\mathcal U_t^{-1}$ is a continuous inverse of
$S_t-\lambda$. Conjugating in the opposite direction proves the
converse, and hence the spectrum equality using
\eqref{eq:cq-source-spectrum}. Apply the polynomial identity
$U_tp(M)=p(B_t)U_t$ to $\Xi$ for the last assertion.
\end{proof}

The heat map also transports the full multiplication, with every
cross term retained. For $f,g\in\cqE$ define
\begin{equation}\label{eq:cq-star}
 f\star_t g=U_t\bigl((U_{-t}f)(U_{-t}g)\bigr)
    =\sum_{n=0}^\infty\frac{(-t/2)^n}{n!}(D^nf)(D^ng).
\end{equation}
We prove the series identity, including convergence. For
$F_t=U_tf$, $G_t=U_tg$, the series
\[
 J_t=\sum_{n=0}^\infty\frac{(-t/2)^n}{n!}
                    (D^nF_t)(D^nG_t)
\]
converges in $\gamma$, uniformly on compact time sets: applying
\eqref{eq:cq-cauchy} with output norm $(p,a/2)$ to both factors
bounds its $n$th term by
\[
 K\left[\frac T2\left(\frac{epa}{2}\right)^{2/p}
                   n^{1-2/p}\right]^n\quad(|t|\leq T),
\]
where $K$ bounds the product of the two input norms
$(p,a/2^p)$, and is finite by the heat estimate. The root test
applies. Cauchy's formula in time permits termwise time derivatives;
continuity of $D$ permits spatial derivatives. Differentiating the
coefficient contributes
$-\tfrac12\sum(-t/2)^n(D^{n+1}F_t)(D^{n+1}G_t)/n!$.
Differentiating the factors contributes the two pure second-derivative
terms with coefficient $-1/4$. The complete product rule therefore
gives $\partial_tJ_t=-D^2J_t/4$ and $J_0=fg$.
Its entire time Taylor coefficients are uniquely
$(-1/4)^kD^{2k}(fg)/k!$, so $J_t=U_t(fg)$.
Replacing the inputs by $U_{-t}f,U_{-t}g$ proves
\eqref{eq:cq-star}. Conjugation of ordinary multiplication proves
associativity, commutativity, unit $1$, and continuity for $\tau_t$,
even when continuity of the series in the source topology is unknown.
The identity $s\star_t f=sf-tf'/2$ recovers exactly $B_t$.
No ideal property of $\cqN$ beyond its stated polynomial invariance
is required for this identity or for the quotient in
Theorem~\ref{thm:cq-moving}.

The transported actual trace image has no common point-evaluation
zero at any nonzero time. Suppose $\lambda$ were such a common
zero. It would annihilate each $U_t(s^n\Xi)$. For every $a\in\cqC$,
the Taylor series of $e^{as}$ converges in $\gamma$ by
Lemma~\ref{thm:cq-density}. Multiplication, $U_t$, and evaluation
are continuous in $\gamma$, so
\[
 0=\sum_{n=0}^\infty\frac{a^n}{n!}U_t(s^n\Xi)(\lambda)
   =e^{a\lambda-ta^2/4}Z_t(\lambda-ta/2).
\]
The last equality follows from the just-proved product formula
and the entire Taylor formula for $Z_t$. At $t\ne0$ varying $a$
forces $Z_t$ to vanish everywhere, a contradiction. This argument
uses only the polynomial generators contained in $N_t$; it assumes
no source-topology convergence of the exponential series. Thus the
source arithmetic spectrum in \eqref{eq:cq-specmove} is carried by
the explicitly transported operator and relation, even though the
transported relation has no common ordinary evaluation zero.

The first two nonconstant polynomial generators make the extra terms
visible without changing any polynomial coefficient:
\begin{align}\label{eq:cq-generators}
 U_t(s\Xi)&=sZ_t-\tfrac t2Z_t',\notag\\
 U_t(s^2\Xi)&=s^2Z_t-t sZ_t'-\tfrac t2Z_t+\tfrac{t^2}4Z_t''.
\end{align}
For $L=\partial_t+\tfrac14D^2$, the complete product rule is
\begin{equation}\label{eq:cq-unit}
 L(AP)=A\bigl(P_t+\tfrac14P_{ss}\bigr)
       +\bigl(A_t+\tfrac14A_{ss}\bigr)P
       +\tfrac12A_sP_s.
\end{equation}
Thus $L(R(s)Z_t)=\tfrac14R''Z_t+\tfrac12R'Z_t'$.
The family $R(B_t)Z_t$ solves the homogeneous heat equation, whereas
the displayed expression is the full defect of replacing it by
$R(s)Z_t$. For a local factorization $Z_t=AP$ with $A\ne0$,
the same formula retains the analytic unit $A$ and its derivatives.
At a simple zero $\lambda(t)$, differentiating
$Z_t(\lambda(t))=0$ gives
\begin{equation}\label{eq:cq-zero}
 \lambda'(t)=\frac14\frac{Z_{ss}}{Z_s}(t,\lambda(t))
 =\frac14\frac{P_{ss}}{P_s}(t,\lambda(t))
    +\frac12\frac{A_s}{A}(t,\lambda(t)).
\end{equation}
This statement is on the open locus $Z_s\ne0$, on which the analytic
implicit-function theorem supplies the local curve; it claims no
chosen zero is simple. At a multiple zero the division is undefined,
and the full equation \eqref{eq:cq-unit} remains the equation to use.

\section{Derivative domains, multivalued operators, and jet extensions}
\label{sec:cq-jets}
Every derivative has the full underlying domain $\cqE$ by
Lemma~\ref{thm:cq-heat}. Continuity in the unnamed source topology is
not assumed. The following gives exact source-quotient operators and
their domains without imposing it. For $k\geq1$ put
\[
 K_k=\{n\in\cqN:D^kn\in\cqN\},\qquad
 W_k=(D^k\cqN+\cqN)/\cqN\subseteq\cqQ.
\]
The derivative correspondence is the actual linear relation
\begin{equation}\label{eq:cq-relation}
 \mathfrak D_k=\{([f],[D^kf]):f\in\cqE\}
                    \subseteq\cqQ\times\cqQ.
\end{equation}
For fixed $[f]$, its output fibre is precisely $[D^kf]+W_k$:
all representatives are $f+n$ with $n\in\cqN$, and linearity gives
all and only the stated outputs. The kernel of its presenting map
$f\mapsto([f],[D^kf])$ is exactly $K_k$. The relation has domain
all of $\cqQ$; its multivalued part at zero is exactly $W_k$.
Thus it is the graph of an operator exactly when $D^k\cqN\subseteq
\cqN$. This is a calculated equivalence, with an explicit obstruction
space, rather than an assumption of invariance.

There is no nonempty domain of quotient classes on which the same
representative formula is unambiguous if $W_k\ne0$ and every
representative of a class is admitted: addition of the same offending
$n$ gives two different outputs for every class. A chosen linear
lift is a different map. More precisely any vector subspace
$A\subseteq\cqE$ supplies a single-valued map on its quotient image
$q(A)$ exactly when $D^k(A\cap\cqN)\subseteq\cqN$; the identical
difference-of-representatives argument proves both directions.
The first heat generator is $-D^2/4$, so its ambiguity space is
$-W_2/4$, with the coefficient and sign unchanged.

\subsection{Exact algebraic spectral defects of the source relation}
\label{sec:cq-sourcealgebra}
These calculations keep the actual $\cqN$ and require only the
already proved $M\cqN\subseteq\cqN$. For every $\lambda\in\cqC$,
division by $x=s-\lambda$ gives the algebraic identities
\begin{equation}\label{eq:cq-source-kercoker}
 \begin{aligned}
 \operatorname{coker}_{\rm alg}(S-\lambda)
     &=\cqE/(\cqN+x\cqE),\\
 \ker(S-\lambda)&\simeq
        (\cqN\cap x\cqE)/(x\cqN),
        & [n]&\longmapsto[n/x].
 \end{aligned}
\end{equation}
Division belongs to the original order class: if $f(\lambda)=0$,
the filled quotient is entire, is bounded on a fixed disk, and
outside $|s|\geq2|\lambda|+1$ its modulus is at most $2|f(s)|$.
Thus $x\cqE$ is exactly the kernel of evaluation at $\lambda$.
The cokernel identity follows from the image of the induced
multiplication. For the second identity, changing $n$ by $xn_0$
changes $n/x$ by $n_0\in\cqN$; its image is in the stated
kernel. Conversely any kernel representative $f$ has
$xf\in\cqN\cap x\cqE$. Finally $[n/x]=0$ exactly when
$n\in x\cqN$. This proves every asserted kernel and image.

The cokernel has only two possibilities. When every $n\in\cqN$
vanishes at $\lambda$, it is one-dimensional, with its exact
identification given by evaluation. Otherwise retain any
$h_\lambda\in\cqN$ with $h_\lambda(\lambda)\ne0$. Then
\begin{equation}\label{eq:cq-source-division}
 \begin{aligned}
 A_\lambda f(s)&=
  \frac{f(s)-f(\lambda)h_\lambda(s)/h_\lambda(\lambda)}{s-\lambda},\\
 f&=\frac{f(\lambda)}{h_\lambda(\lambda)}h_\lambda
                                      +(s-\lambda)A_\lambda f
 \end{aligned}
\end{equation}
proves surjectivity of $S-\lambda$, and its exact remaining kernel is
\begin{equation}\label{eq:cq-source-splitting}
 \ker(S-\lambda)\simeq
 \cqN/\bigl((s-\lambda)\cqN\oplus\cqC h_\lambda\bigr),
 \qquad n\longmapsto[A_\lambda n].
\end{equation}
The map is onto because a kernel class $[g]$ is the image of
$n=(s-\lambda)g\in\cqN$. Its kernel consists exactly of
$n=(s-\lambda)n_0+c h_\lambda$, $n_0\in\cqN$, by
\eqref{eq:cq-source-division}. The sum is direct by evaluation.
In this second case algebraic invertibility is therefore equivalent
to the exact identity
$\cqN=(s-\lambda)\cqN\oplus\cqC h_\lambda$.
On that locus the representative formula $[f]\mapsto[A_\lambda f]$
is independent of representatives, since $A_\lambda\cqN\subseteq
\cqN$; multiplication in \eqref{eq:cq-source-division} proves the
right inverse, and $A_\lambda((s-\lambda)f)=f$ proves the left
inverse. These are algebraic statements; continuity of this inverse
in $\tau$ has not been inserted.

The first derivative and the heat generator have exactly the same
frozen invariance question for the source subspace:
\begin{equation}\label{eq:cq-source-D-equivalence}
 D\cqN\subseteq\cqN
 \quad\Longleftrightarrow\quad D^2\cqN\subseteq\cqN
 \quad\Longleftrightarrow\quad -\tfrac14D^2\cqN\subseteq\cqN.
\end{equation}
Indeed the product rule is
$D^2(sn)-sD^2n=2Dn$. Since $sn\in\cqN$, second-derivative
invariance puts both terms on the left in $\cqN$, proving
first-derivative invariance. Iteration proves the reverse implication;
the retained nonzero scalar $-1/4$ proves the last equivalence.
If these equivalent inclusions hold, all $D^j\Xi$ belong to
$\cqN$. At each $\lambda$ at least one is nonzero: otherwise
the Taylor series would make $\Xi$ identically zero. Thus every
$S-\lambda$ is then algebraically surjective by
\eqref{eq:cq-source-division}. Its prescribed topological spectral
points would be represented entirely by the kernel
\eqref{eq:cq-source-splitting} or a noncontinuous algebraic inverse.
This computes the precise alternatives without assigning either
one to the source's stated spectral set.

There is also a full multiplicity consequence of algebraic derivative
descent. Writing $d[f]=[Df]$ on its exact descent locus, the product
rule gives $[d,S]=I$. For $0\ne v\in\ker(S-\lambda)$,
\begin{equation}\label{eq:cq-source-weylladder}
 (S-\lambda)d^kv=-k d^{k-1}v,
 \qquad (S-\lambda)^kd^kv=(-1)^k k!v\quad(k\geq0),
\end{equation}
with zero right side in the first identity at $k=0$.
To prove induction, commute the next $d$ using
$(S-\lambda)d=d(S-\lambda)-I$; the coefficient $-k$ becomes
$-(k+1)$. Iteration proves the second identity. A finite linear
dependence among $v,dv,\ldots,d^nv$, with largest nonzero
coefficient at index $j$, becomes that coefficient times
$(-1)^j j!v$ after applying $(S-\lambda)^j$, a contradiction.
Thus every such eigenvector generates an infinite independent
generalized eigenvector sequence. No finite source multiplicity
assertion has been assumed.

Finally even a single forward nonzero heat inclusion has a calculated
consequence. Fix $t\ne0$ and $c=t/2$, and define
\begin{equation}\label{eq:cq-source-heatcommutator}
 \begin{gathered}
 A_0=U_t,\qquad A_{k+1}=MA_k-A_kM,\qquad A_k=P_k(D)U_t,\\
 P_0(X)=1,\qquad P_{k+1}(X)=cX P_k(X)-P_k'(X).
 \end{gathered}
\end{equation}
The identity follows by induction from $[M,U_t]=cDU_t$ and
$[M,D^j]=-jD^{j-1}$: commuting $M$ through $P_k(D)U_t$
contributes $-P_k'(D)U_t+cP_k(D)DU_t$.
The degree of $P_k$ is exactly $k$, with leading coefficient
$c^k\ne0$. Under $U_t\cqN\subseteq\cqN$, induction with
$M\cqN\subseteq\cqN$ gives $A_k\cqN\subseteq\cqN$;
subtracting the lower-degree terms of $P_k$ then gives exactly
\begin{equation}\label{eq:cq-source-heatconsequence}
 D^kU_t\cqN\subseteq\cqN\quad(k\geq0).
\end{equation}
In particular every derivative of the nonzero $Z_t=U_t\Xi$ lies
in $\cqN$. At every point one derivative is nonzero by the entire
Taylor theorem, so every $S-\lambda$ is algebraically surjective
in this case too. This proof uses only the one forward inclusion.
For the stronger equality $U_t\cqN=\cqN$, it also proves
$D\cqN\subseteq\cqN$, by writing each $n\in\cqN$ as
$U_tm$ in \eqref{eq:cq-source-heatconsequence}. None of these
exact implications presumes the inclusion or equality for the actual
source closure. The explicit obstruction spaces above remain the
calculated source question.


Even when $W_k$ is not closed, the smallest source-closed enlargement
of the target relation subspace is explicit:
\begin{equation}\label{eq:cq-target}
 [f]\longmapsto[D^kf]:
 \cqQ\longrightarrow
 \cqE/\overline{\cqN+D^k\cqN}^{\,\tau}.
\end{equation}
It is well defined algebraically. Equip its source with the quotient
of the graph topology $f\mapsto(f,D^kf)$ in
$(\cqE,\tau)^2$ to obtain a continuous map. This specifies a
topology change explicitly, without claiming it is the original one.
The lost part of the original target $\cqQ$ is exactly
$\overline{\cqN+D^k\cqN}^{\,\tau}/\cqN$.

Here is an extension which retains correlated derivatives rather than
quotienting them away. For $j\geq1$ define the bijection
\[
 T_j(f_0,\ldots,f_j)
   =(f_0,f_1-Df_0,\ldots,f_j-D^jf_0).
\]
Its inverse is $(f_0,g_1,\ldots,g_j)\mapsto
(f_0,g_1+Df_0,\ldots,g_j+D^jf_0)$.
Give $V_j=\cqE^{j+1}$ the topology transported by $T_j$ from
$(\cqE,\tau)\times(\cqE,\gamma)^j$. The residual factors
are explicitly given the proved Hausdorff growth topology; the
source factor keeps its original topology. For $J_jf=(f,Df,\ldots,D^jf)$,
$T_j(J_j\cqN)=\cqN\oplus0^j$, so $J_j\cqN$ is closed.
Define
\begin{equation}\label{eq:cq-jetextension}
 X_j=V_j/J_j\cqN.
\end{equation}
Then the map
\[
 [f_0,\ldots,f_j]\longmapsto
 ([f_0],f_1-Df_0,\ldots,f_j-D^jf_0)
\]
is a homeomorphism $X_j\simeq\cqQ\oplus(\cqE,\gamma)^j$; independence,
the inverse, and continuity follow directly from $T_j$ and the
quotient definition. In particular there is the split exact sequence
\begin{equation}\label{eq:cq-jetseq}
 0\longrightarrow\cqE^j\longrightarrow X_j
 \xrightarrow{[f_0,\ldots,f_j]\mapsto[f_0]}\cqQ
 \longrightarrow0,
 \qquad [f]\longmapsto[J_jf].
\end{equation}
Its relation on a source generator retains every product term:
\begin{equation}\label{eq:cq-rawjet}
 D^k(R\Xi)=\sum_{h=0}^k\binom kh R^{(h)}\Xi^{(k-h)}
                    \quad(0\leq k\leq j).
\end{equation}
The binomial formula follows by induction from the ordinary product
rule, including the two edge terms at each step.

The original spectral coordinate acts on $V_j$ by
\[
 (\mathcal M_j f)_k=s f_k+k f_{k-1}\quad(0\leq k\leq j),
 \qquad f_{-1}=0.
\]
The identity $D^k(sf)=sD^kf+kD^{k-1}f$ proves
$\mathcal M_jJ_jf=J_j(sf)$; it therefore preserves the exact
relation subspace. In the $T_j$ coordinates it acts by $S$ on
$\cqQ$ and by $g_k\mapsto sg_k+kg_{k-1}$ on the residual terms
($g_0=0$). This also proves continuity in the stated topology.
Notice that \eqref{eq:cq-jetseq} adds a free residual module:
the spectrum of the entire extension is not claimed to equal the
source spectrum. The arithmetic quotient and its split inclusion
are explicitly preserved.

\section{A fully computed analytic realization and its exact comparison}
\label{sec:cq-analytic}
We now compute the closures and the failure of descent in specified
topologies. This does not identify them with $\tau$.
Let $\kappa$ be the topology of uniform convergence on compact sets
on the same set $\cqE$, and define
\[
 N_\gamma=\overline{\cqC[s]\Xi}^{\,\gamma},\qquad
 N_\kappa=\overline{\cqC[s]\Xi}^{\,\kappa}.
\]
Evaluations of all derivatives are continuous in both topologies by
Cauchy's formula. The growth topology is finer than $\kappa$, so
$N_\gamma\subseteq N_\kappa$.
If $\lambda$ is a zero of $\Xi$ with its actual multiplicity $m_\lambda$,
every element of either closed subspace has vanishing derivatives
of orders $0,\ldots,m_\lambda-1$ at $\lambda$.

\begin{cqlemma}[Entire division with its growth and topology retained]
\label{thm:cq-division}
Let $g\in\cqE$ with $g(0)\ne0$. If $f\in\cqE$ and $h=f/g$
extends to an entire function, then $h\in\cqE$. For $1<p<2$,
$a>0$, put $b=a/(6\cdot2^p)$. There is the linear division estimate
\begin{equation}\label{eq:cq-division-bound}
 \|f/g\|_{p,a}\leq
 \left(\frac{\max(1,\|g\|_{p,b})}{|g(0)|}\right)^3
 \|f\|_{p,b}.
\end{equation}
Multiplication by $g$ is a topological linear isomorphism from
$(\cqE,\gamma)$ onto the closed subspace of functions divisible by
$g$, with its subspace topology. Compact-topology division is also
continuous.
\end{cqlemma}
\begin{proof}
Choose $R>0$ with no zero of $g$ on $|s|=R$. Factor its finitely
many zeros in the disk, retaining multiplicities. The residual is
holomorphic and zero-free there, so its log modulus is harmonic.
For $|z|<R$ the boundary mean of $\log|Re^{i\theta}-z|$ equals
$\log R$: expand $\log(1-(z/R)e^{-i\theta})$ in its uniformly
convergent series and take real parts. The mean-value formula for
the residual gives the exact Jensen identity
\[
 \frac1{2\pi}\int_0^{2\pi}\log|g(Re^{i\theta})|\,d\theta
 =\log|g(0)|+\sum_{|z|<R}m_z\log(R/|z|)\geq\log|g(0)|.
\]
There is no zero at zero. With
$\log^+x=\max(0,\log x)$ and $\log^-x=\max(0,-\log x)$,
the inequality $\log^+|h|\leq\log^+|f|+\log^-|g|$ gives
\begin{equation}\label{eq:cq-division-mean}
 \frac1{2\pi}\int_0^{2\pi}\log^+|h(Re^{i\theta})|\,d\theta
 \leq\log^+M_f(R)+\log^+M_g(R)-\log|g(0)|.
\end{equation}
The right side is nonnegative since $M_g(R)\geq|g(0)|$.
The harmonic Poisson extension of the continuous nonnegative boundary
function $\log^+|h|$ majorizes $\log|h|$ by the maximum principle
and is nonnegative; hence it majorizes $\log^+|h|$ in the disk.
For $|s|\leq r<R$ the Poisson kernel is at most $(R+r)/(R-r)$.
Take $R\downarrow2r$ through circles avoiding zeros. Continuity of
the two maximum functions gives, for every $r>0$,
\begin{equation}\label{eq:cq-division-radial}
 \log^+M_h(r)\leq
 3\bigl(\log^+M_f(2r)+\log^+M_g(2r)-\log|g(0)|\bigr).
\end{equation}
No pointwise lower bound for $g$ on a large circle has been assumed.
For $f\ne0$, put $F=\|f\|_{p,b}>0$ and apply this to $f/F,h/F$.
Then $\log^+M_{f/F}(2r)\leq b(2r)^p$, and
$\log^+M_g(2r)\leq\log\max(1,\|g\|_{p,b})+b(2r)^p$.
The coefficient of $r^p$ is $6b2^p=a$, giving
\eqref{eq:cq-division-bound}; continuity fills in $r=0$.
The zero numerator is immediate. All $p,a$ are allowed, so $h$
belongs to the original order class and division is continuous.
Multiplication is continuous by Lemma~\ref{thm:cq-density}.
The requisite vanishing jets at every zero are closed conditions;
local Taylor division identifies them with entire divisibility.
The estimate proves this closed set is exactly $g\cqE$.
For the compact topology choose $R>r$ with no boundary zero and
$\delta=\min_{|s|=R}|g(s)|>0$. The maximum principle for the filled
entire quotient gives $M_h(r)\leq\delta^{-1}M_f(R)$, proving
compact continuity too.
\end{proof}

For general nonzero $g\in\cqE$, choose and retain a point $c$ with
$g(c)\ne0$ and use disks centered at $c$. The functions and spectral
coordinate remain unchanged. The bound
$|s+c|^p\leq2^{p-1}(|s|^p+|c|^p)$ proves continuity of translations
by $c$ and $-c$ in $\gamma$; compact containment proves it in
$\kappa$. Thus all conclusions hold with the stated translated
estimates. The original center $c=0$ already applies to $\Xi$.

\begin{cqproposition}[The computed closures are the same subspace]
\label{thm:cq-equalclosures}
With their definitions unchanged,
\begin{equation}\label{eq:cq-equalclosures}
 \begin{aligned}
 N_\gamma=N_\kappa&=\Xi\cqE\\
 &=\{f\in\cqE:D^hf(\lambda)=0\text{ for every }\Xi(\lambda)=0,
                                      \ 0\leq h<m_\lambda\}.
 \end{aligned}
\end{equation}
This is equality of subsets of the original ring, not an
identification of $\gamma$, $\kappa$, or $\tau$.
\end{cqproposition}
\begin{proof}
Jet continuity puts both closures inside the displayed vanishing
set and gives $N_\gamma\subseteq N_\kappa$. Local Taylor division
makes $h=f/\Xi$ entire for any $f$ in that set; the division lemma
then puts $h$ in $\cqE$. Its Taylor polynomials $p_n$ converge in
$\gamma$, and multiplication gives $p_n\Xi\to f$ there. Thus the
vanishing set is contained in $N_\gamma$, proving all equalities.
The algebraic-ideal equality is the conclusion of this argument;
it was not substituted for an unevaluated closure.
\end{proof}

\begin{cqtheorem}[Compact closure, multiplicities, and coordinate spectrum]
\label{thm:cq-compact}
The compact closure is exactly
\begin{equation}\label{eq:cq-compactideal}
 \begin{aligned}
 N_\kappa
 &=\{f\in\cqE:f/\Xi\text{ extends to an entire function}\}\\
 &=\{f\in\cqE:D^hf(\lambda)=0\ (h<m_\lambda),\ \Xi(\lambda)=0\}.
 \end{aligned}
\end{equation}
On $Q_\kappa=(\cqE,\kappa)/N_\kappa$, multiplication by $s$ has
spectrum exactly the zeros of $\Xi$ in the original $s$ coordinate.
For a zero $\lambda$ of multiplicity $m$, there is a continuous
surjection
\[
 Q_\kappa\longrightarrow\cqC[\epsilon]/(\epsilon^m),\qquad
 [f]\longmapsto\sum_{h=0}^{m-1}\frac{D^hf(\lambda)}{h!}\epsilon^h,
\]
intertwining multiplication by $s$ with $\lambda+\epsilon$.
\end{cqtheorem}
\begin{proof}
The vanishing conditions are closed and hold on every $p\Xi$,
so the closure is contained in the last displayed set. Local Taylor
division at each isolated zero identifies that set with entire
divisibility by $\Xi$. For such an $f$, put $g=f/\Xi\in\cqO(\cqC)$.
Its Taylor polynomials $p_n$ converge uniformly on every compact;
$p_n\Xi\to f$ in $\kappa$. Every $p_n\Xi$ belongs to the original
order class even though no growth assertion about $g$ is needed.
This proves \eqref{eq:cq-compactideal} with the closure retained.
The local-jet map is independent of representatives; the polynomials
$1,(s-\lambda),\ldots,(s-\lambda)^{m-1}$ prove surjectivity.
Leibniz's rule gives the stated action, and the finite jet evaluations
are continuous.

If $\Xi(\lambda)\ne0$, define
\[
 T_\lambda f(s)=
 \frac{f(s)-f(\lambda)\Xi(s)/\Xi(\lambda)}{s-\lambda},
\]
filling in the removable value. Division by a linear factor preserves
the order bound: outside a disk about $\lambda$ the denominator is
bounded below, and inside a disk the filled entire function is bounded.
It is continuous in $\kappa$ by Cauchy's formula on a larger disk
applied to the numerator. For $n\in N_\kappa$ the quotient
$T_\lambda n/\Xi=(n/\Xi-(n/\Xi)(\lambda))/(s-\lambda)$ is entire;
thus it descends. Direct multiplication proves
$(S_\kappa-\lambda)[T_\lambda f]=[f]$.
Applying $T_\lambda$ to $(s-\lambda)f$ gives $f$ exactly, proving
the other inverse identity. If $\Xi(\lambda)=0$, evaluation at
$\lambda$ descends to a nonzero continuous functional which annihilates
the image of $S_\kappa-\lambda$. Hence that operator is not
surjective. This proves the spectrum claim.
\end{proof}

The same coordinate spectrum and the same finite-jet surjections hold
for $Q_\gamma=(\cqE,\gamma)/\Xi\cqE$. Indeed the numerator in
$T_\lambda f$ depends continuously on $f$ in $\gamma$. Division
by $s-\lambda$, on its entire-divisibility domain, is continuous by
Lemma~\ref{thm:cq-division} with a nonzero center retained when
needed. The already proved two inverse identities therefore give a
continuous resolvent in $\gamma$ as well. The evaluation obstruction
at each zero is continuous in that topology too. This proves the
claim, without transferring a resolvent through an unproved topology
identification.

The ambient topologies $\kappa$ and $\gamma$ are in fact distinct.
The original functions $f_n(s)=\exp(ns-n^2)$ tend to zero uniformly
on each compact set. On the positive real axis the maximum of
$nx-x^{3/2}$ is $4n^3/27$, attained at $x=4n^2/9$. Since
$\operatorname{Re}s\leq|s|$, that axis attains the full weighted
supremum, giving the exact value
\[
 \|f_n\|_{3/2,1}=\exp(4n^3/27-n^2)\longrightarrow\infty.
\]
Thus equality of the closed arithmetic relation subspaces and of
their spectral sets does not identify the ambient topologies. The
identity induces a continuous linear bijection $Q_\gamma\to Q_\kappa$;
the next theorem proves that its inverse is not continuous.


\subsection{The actual difference between the two analytic quotient topologies}
\label{sec:cq-quotient-topologies}
Retain the original entire-function set $\cqE=\mathcal H_{\leq1}$,
the original coordinate $s$, and
$\Xi(s)=\xi(\tfrac12+is)$ from \eqref{eq:cq-Xi}.
Let $\Lambda$ be the set of its distinct zeros and let
$m_\lambda=\operatorname{ord}_\lambda\Xi\geq1$ be each actual
multiplicity. Put
\[
 I=\Xi\cqE,\qquad
 Q_\kappa=(\cqE,\kappa)/I,\qquad
 Q_\gamma=(\cqE,\gamma)/I,\qquad
 \mathcal P_\Xi=\prod_{\lambda\in\Lambda}\cqC^{m_\lambda}.
\]
The product has its ordinary product topology. Coordinates in
$\cqC^{m_\lambda}$ are the actual derivatives of orders
$0,\ldots,m_\lambda-1$, with no factorial rescaling. Define
\begin{equation}\label{eq:cq-total-jet-map}
 J_\Xi f=(f^{(j)}(\lambda))_{
          \lambda\in\Lambda,\ 0\leq j<m_\lambda},
 \qquad \mathcal V_\Xi=J_\Xi(\cqE)\subset\mathcal P_\Xi.
\end{equation}
Every point of $\mathcal V_\Xi$ retains the entire input's arithmetic
jets. The restriction to this image is essential.

\begin{cqtheorem}[Jet topology and its completion]
\label{thm:cq-jet-topology}
The map induced by \eqref{eq:cq-total-jet-map} is a topological
linear isomorphism
\[
 \overline J_\Xi:Q_\kappa\xrightarrow{\ \cong\ }
 \mathcal V_\Xi
\]
onto the subspace topology from $\mathcal P_\Xi$.
The image is dense and proper. Consequently $Q_\kappa$ is
incomplete, and its Hausdorff completion is canonically
$\mathcal P_\Xi$ via this map. In contrast $Q_\gamma$ is a
Fr\'echet space and is complete. The identity on the same
underlying quotient vector space induces a continuous linear
bijection
\begin{equation}\label{eq:cq-topology-comparison-strict}
 \iota:Q_\gamma\longrightarrow Q_\kappa
\end{equation}
whose inverse is not continuous.
\end{cqtheorem}

\begin{proof}
We first construct each finite interpolation map with its constants.
For a finite nonempty set $F\subset\Lambda$, write
$d_F=\sum_{\lambda\in F}m_\lambda$ and set
\[
 B_{\lambda,F}(s)=
 \prod_{\substack{\nu\in F\\\nu\ne\lambda}}
 (s-\nu)^{m_\nu}.
\]
The empty product is $1$. Since
$B_{\lambda,F}(\lambda)\ne0$, the germ
$(s-\lambda)^j/(j!B_{\lambda,F}(s))$ is holomorphic at
$\lambda$. Denote its Taylor polynomial through degree
$m_\lambda-1$, in powers of the unchanged difference
$s-\lambda$, by $T_{\lambda,m_\lambda-1}$ of that germ. Put
\begin{equation}\label{eq:cq-Hermite-basis-unscaled}
 L_{\lambda,j,F}(s)=B_{\lambda,F}(s)
 T_{\lambda,m_\lambda-1}
 \left(\frac{(s-\lambda)^j}{j!B_{\lambda,F}(s)}\right),
 \qquad 0\leq j<m_\lambda.
\end{equation}
This is a polynomial of degree at most $d_F-1$.
For $\nu\in F\setminus\{\lambda\}$ it vanishes to order
at least $m_\nu$ at $\nu$. At $\lambda$ its difference from
$(s-\lambda)^j/j!$ vanishes to order at least $m_\lambda$.
Therefore, for all $\nu\in F$ and $0\leq k<m_\nu$,
\[
 L_{\lambda,j,F}^{(k)}(\nu)
   =\begin{cases}1,&(\nu,k)=(\lambda,j),\\0,&\text{otherwise}.
     \end{cases}
\]
It follows that the linear map to polynomials
\begin{equation}\label{eq:cq-finite-Hermite-lift}
 H_F(v)(s)=\sum_{\lambda\in F}\sum_{j=0}^{m_\lambda-1}
              v_{\lambda,j}L_{\lambda,j,F}(s)
\end{equation}
realizes every prescribed finite array of actual derivatives
$v=(v_{\lambda,j})$. For each $R>0$ define
\[
 M_f(R)=\max_{|s|\leq R}|f(s)|,\qquad
 C_{R,F}=\sum_{\lambda\in F}\sum_{j=0}^{m_\lambda-1}
                  M_{L_{\lambda,j,F}}(R).
\]
The triangle inequality gives the exact finite bound
\begin{equation}\label{eq:cq-Hermite-lift-bound}
 M_{H_F(v)}(R)\leq C_{R,F}
             \max_{\lambda\in F,\,j<m_\lambda}|v_{\lambda,j}|.
\end{equation}
When $F$ is empty set $H_F=0$, $C_{R,F}=0$ and the maximum
over that empty set to $0$.

The kernel of $J_\Xi$ is precisely $I$. Indeed $J_\Xi f=0$
means that $f/\Xi$ has a removable singularity at each zero of
$\Xi$ by its local Taylor expansions. Away from those zeros it
is already holomorphic. Thus it extends to an entire function,
which belongs to the original order class $\cqE$ by
Lemma~\ref{thm:cq-division}; consequently $f\in\Xi\cqE$.
The converse follows by the product rule at each zero. This proves
that $\overline J_\Xi$ is a well-defined linear bijection onto
$\mathcal V_\Xi$, with inverse sending the jet array of any
$f\in\cqE$ to its class $[f]$. The kernel calculation proves
that this inverse is independent of the choice of such $f$.

The quotient topology on $Q_\kappa$ is generated by
\begin{equation}\label{eq:cq-compact-quotient-seminorm}
 q_R([f])=\inf_{h\in\cqE}M_{f-\Xi h}(R),\qquad R>0.
\end{equation}
For completeness, these are seminorms because they are infima of
the representatives' seminorms modulo a linear subspace. The
identity
$\{q_R<\epsilon\}=\pi_\kappa(\{M_f(R)<\epsilon\})$
holds for every $\epsilon>0$, and the quotient map
$\pi_\kappa$ is open since the inverse image of the image of
an open set $U$ is the union $U+I$ of its translates.
The increasing disk seminorms generate the compact-open topology;
the displayed identity therefore proves the assertion about the
quotient topology. Integer $R\geq1$ suffice.

Fix $R>0$. Choose $R'>R$ such that the circle $|s|=R'$
contains no zero of $\Xi$, and put
$F=\{\lambda\in\Lambda:|\lambda|<R'\}$, a finite set.
For $f\in\cqE$ form
$H=H_F((f^{(j)}(\lambda))_{\lambda\in F,j<m_\lambda})$.
The function $(f-H)/\Xi$ is holomorphic throughout $|s|<R'$
by the exact finite interpolation just proved. Its Taylor
polynomials $p_n(s)$ at $s=0$ consequently converge uniformly on
$|s|\leq R$ to $(f-H)/\Xi$. In particular $p_n\in\cqE$,
and
\[
 M_{f-\Xi p_n}(R)
 \leq M_H(R)+M_\Xi(R)
        M_{(f-H)/\Xi-p_n}(R).
\]
Taking the infimum over admissible representatives and then the
limit gives
\begin{equation}\label{eq:cq-quotient-from-finite-jets}
 q_R([f])\leq M_H(R)
 \leq C_{R,F}
       \max_{\lambda\in F,\,j<m_\lambda}|f^{(j)}(\lambda)|.
\end{equation}
No entire extension or growth bound for $(f-H)/\Xi$ outside this
disk was used. In particular there is no additional division
assumption hidden in the approximation. If $F$ is empty, this
proof gives $q_R=0$, which is compatible with the other quotient
seminorms separating points.

Conversely, fix a finite set $F\subset\Lambda$, and choose
$R>\max_{\lambda\in F}|\lambda|$ and
$0<\delta<R-\max_{\lambda\in F}|\lambda|$.
Cauchy's formula on the radius-$\delta$ circle about each
$\lambda$ gives, for every $h\in\cqE$,
\[
 |f^{(j)}(\lambda)|
 =|(f-\Xi h)^{(j)}(\lambda)|
 \leq j!\delta^{-j}M_{f-\Xi h}(R),
 \qquad\lambda\in F,\quad j<m_\lambda.
\]
Here $0!=1$, so this bound includes the zeroth derivative without
altering its value. Taking infima gives
\begin{equation}\label{eq:cq-finite-jets-from-quotient}
 \max_{\lambda\in F,\,j<m_\lambda}|f^{(j)}(\lambda)|
 \leq\left(\max_{\lambda\in F,\,j<m_\lambda}
                        j!\delta^{-j}\right)q_R([f]).
\end{equation}
Equations \eqref{eq:cq-quotient-from-finite-jets} and
\eqref{eq:cq-finite-jets-from-quotient} prove continuity in both
directions, with the exact stated topologies.

We next verify that the zero set needed for incompleteness is
indeed infinite and has an enumeration escaping every compact
set. The source Hadamard product for the original even $\Xi$,
retaining its nonzero constant $\Xi(0)$, implies that finitely
many zeros would make $\Xi$ a polynomial. This is impossible:
at the original coordinate $s=-i(2n-\tfrac12)$, $n\geq1$,
the defining completion gives exactly
\begin{equation}\label{eq:cq-Xi-nonpolynomial-values}
 \Xi\bigl(-i(2n-\tfrac12)\bigr)
 =\xi(2n)=n(2n-1)(n-1)!\,\pi^{-n}\zeta(2n).
\end{equation}
For integer $n\geq2$, $\zeta(2n)>1$ by its positive Dirichlet
series, and the last $\lfloor n/2\rfloor$ factors of $(n-1)!$
are at least $n/2$. Thus the logarithm of the right hand side is
at least
\[
 \log(n(2n-1))+\lfloor n/2\rfloor\log(n/2)-n\log\pi,
\]
which exceeds $d\log(2n)$ for every fixed $d$ and all
sufficiently large $n$. Polynomial values of degree $d$ on these
points are bounded by a constant times $(2n)^d$, a contradiction.
Since $\Xi$ is nonzero and entire, its zeros are isolated, each
multiplicity is finite and each closed disk contains finitely
many distinct zeros. We may therefore enumerate all distinct
zeros as $\Lambda=\{\lambda_1,\lambda_2,\ldots\}$ with
$|\lambda_n|\to\infty$, retaining each $m_{\lambda_n}$.
This argument uses neither reality nor simplicity of any zero.

Every projection of $\mathcal V_\Xi$ onto finitely many full
zero-jet blocks is onto, by \eqref{eq:cq-finite-Hermite-lift}.
Such finite coordinate cylinders form a base of the product
topology. Consequently $\mathcal V_\Xi$ is dense in
$\mathcal P_\Xi$. It is proper: define the particular array
$v\in\mathcal P_\Xi$ by
\begin{equation}\label{eq:cq-nonrealizable-jet-array}
 v_{\lambda_n,0}=\exp(|\lambda_n|^2),\qquad
 v_{\lambda_n,j}=0\quad(1\leq j<m_{\lambda_n}).
\end{equation}
For each $f\in\cqE$ there is a finite constant $A_f>0$ with
$M_f(r)\leq A_f\exp(r^{3/2})$ for every $r\geq0$; this is
exactly the order-at-most-one bound, extended over a bounded
interval by increasing $A_f$. Were $J_\Xi f=v$, we would have
\[
 \exp(|\lambda_n|^2)
 \leq A_f\exp(|\lambda_n|^{3/2})\quad\text{for all }n,
\]
which contradicts $|\lambda_n|\to\infty$.

The incompleteness can be exhibited by an actual sequence of
polynomial classes. Set $F_n=\{\lambda_1,\ldots,\lambda_n\}$,
$H_n=H_{F_n}(v|_{F_n})$, and $x_n=[H_n]\in Q_\kappa$.
For each fixed zero-jet block, $J_\Xi H_n$ agrees with $v$
from some index onward; hence $J_\Xi H_n\to v$ in the product.
Equivalently, for each fixed $R$, the finite set in
\eqref{eq:cq-quotient-from-finite-jets} lies in every sufficiently
large $F_n$, so that inequality gives
\[
 q_R(x_n-x_k)=0\qquad(n,k\text{ sufficiently large}).
\]
Thus $(x_n)$ is Cauchy in $Q_\kappa$. If it converged to a
class $[f]$, continuity of $\overline J_\Xi$ and uniqueness of
limits in the product would give $J_\Xi f=v$, already disproved.
For example the product is metrized by the complete
translation-invariant metric
\[
 d(v,w)=\sum_{n=1}^\infty2^{-n}
       \min\!\left(1,\max_{0\leq j<m_{\lambda_n}}
                          |v_{\lambda_n,j}-w_{\lambda_n,j}|\right).
\]
A Cauchy sequence for this metric is Cauchy in each finite
coordinate block, so has a coordinatewise limit; controlling a
finite initial sum and the geometric tail proves convergence in
the metric. This proves product completeness directly. The
isomorphism $\overline J_\Xi$, being linear and continuous
with a continuous linear inverse, preserves the additive
uniformities. Its dense embedding in this complete product
therefore gives the claimed Hausdorff completion.

Finally, we prove the completeness assertion for the stronger
quotient explicitly. The growth seminorms are
\[
 \|f\|_{p,a}=\sup_{s\in\cqC}|f(s)|e^{-a|s|^p},
 \qquad 1<p<2,\quad a>0.
\]
The countable family $p=1+1/j$, $a=1/k$ for integers $j\geq2$
and $k\geq1$ defines the same topology: for arbitrary $1<p<2,a>0$
choose $p_j<p$, and the finite maximum of
$r^{p_j}-ar^p$ controls $\|f\|_{p,a}$ by a constant times
$\|f\|_{p_j,1}$. Each countable seminorm belongs to the original
family, proving both comparisons. Enumerate this countable family
as $(u_\ell)_{\ell\geq1}$ and put
$P_n=\max_{1\leq\ell\leq n}u_\ell$.
If $(f_n)$ is Cauchy in all these seminorms, it is Cauchy uniformly
on each compact set, and has an entire limit $f$. In the weighted
Cauchy inequality for any fixed seminorm, passage to the
pointwise limit in the second index yields the same uniform
weighted bound for $f_n-f$. Thus $f$ has all the finite growth
seminorms, lies in the original $\cqE$, and $f_n\to f$ in
$\gamma$. This proves completeness of $(\cqE,\gamma)$.

The ideal $I$ is closed: the proved kernel description is an
intersection of kernels of derivative evaluations, and Cauchy's
estimate shows that each of these evaluations is continuous for
$\gamma$. The quotient is therefore Hausdorff. Its topology has
the countable increasing seminorm family
\[
 \widetilde P_n([f])=\inf_{h\in I}P_n(f+h).
\]
Let $(y_n)$ be a Cauchy sequence in this quotient. Choose a
strictly increasing subsequence of indices $n_k$ so far out that
\[
 \widetilde P_k(y_{n_{k+1}}-y_{n_k})<2^{-k-1}
 \quad(k\geq1).
\]
This is possible by the Cauchy property and by taking $n_k$ past
the threshold for the $k$th seminorm before selecting the next
index. Choose a representative $d_k\in\cqE$ of each displayed
difference with $P_k(d_k)<2^{-k}$. For any fixed $j$, the tail
$k\geq j$ satisfies $P_j(d_k)\leq P_k(d_k)<2^{-k}$.
Consequently the partial sums of $\sum_{k\geq1}d_k$ are
Cauchy in $\gamma$ and converge in $\cqE$. Adding a
representative of $y_{n_1}$ gives a function whose quotient class
is the limit of $(y_{n_k})$. The full Cauchy sequence $(y_n)$
then has the same limit: for each quotient neighbourhood choose
a smaller additive neighbourhood whose sum with itself is
contained in it, and combine the Cauchy tail with a convergent
subsequence term. The countable seminorm topology is metrizable,
so this proves quotient completeness and the Fr\'echet assertion.

The identity $(\cqE,\gamma)\to(\cqE,\kappa)$ is continuous:
for every $R,p,a$, one has $M_f(R)\leq e^{aR^p}\|f\|_{p,a}$.
It preserves the identical algebraic subspace $I$, so the
quotient universal property proves continuity and bijectivity
of \eqref{eq:cq-topology-comparison-strict}. If its inverse
were continuous, it would be a continuous linear map, hence
uniformly continuous for the additive uniformities. The Cauchy
sequence $(x_n)$ constructed above would then be Cauchy in the
complete $Q_\gamma$, would converge there, and by continuity of
$\iota$ would converge in $Q_\kappa$, a contradiction. This
proves that its inverse is not continuous. No invariance of
metric completeness under arbitrary nonlinear homeomorphisms
has been used.
\end{proof}

The original coordinate action remains visible after completion.
For $v=J_\Xi f$, the exact product rule gives
\begin{equation}\label{eq:cq-original-s-on-completed-jets}
 \bigl(J_\Xi(sf)\bigr)_{\lambda,j}
 =\lambda v_{\lambda,j}+jv_{\lambda,j-1},
 \qquad 0\leq j<m_\lambda,
\end{equation}
where the second term is $0$ when $j=0$. Each output coordinate
depends on finitely many input coordinates, so this extends to
a continuous linear operator on the entire product completion.
The factors $j$ are retained because these coordinates are actual
derivatives. The completion adds precisely the nonrealizable
jet arrays, including \eqref{eq:cq-nonrealizable-jet-array}; it
does not add an entire function of order at most one representing
such an array. Neither analytic quotient topology has been
identified with the source topology $\tau$.


\subsection{Exact derivative domains inside the analytic relation}
Let $I=\Xi\cqE=N_\gamma=N_\kappa$. For $k\geq1$ the exact
representative domain in the relation space is
\begin{equation}\label{eq:cq-Ik}
 \begin{aligned}
 I\cap D^{-k}I=\{\Xi h:\ &h\in\cqE,\quad
 [z^n](A_\lambda(\lambda+z)h(\lambda+z))=0\\
 &\text{for }\max(0,k-m_\lambda)\leq n<k,
                   \text{ at every zero }\lambda\}.
 \end{aligned}
\end{equation}
where the original local unit is
$A_\lambda(s)=\Xi(s)/(s-\lambda)^{m_\lambda}$.
To prove this, write $\Xi h=z^m\sum_{n\geq0}b_nz^n$ at
$\lambda$, with $m=m_\lambda$ and
$b_n=[z^n](A_\lambda h)$. Its $k$th derivative is
\[
 \sum_{m+n\geq k}\frac{(m+n)!}{(m+n-k)!}
               b_nz^{m+n-k}.
\]
Every factorial coefficient is nonzero. Vanishing to order at least
$m$ is equivalent to $b_n=0$ exactly for
$\max(0,k-m)\leq n<k$, as stated. The derivative has the original
order bound, so the vanishing condition is equivalent to membership
in $I$ by \eqref{eq:cq-equalclosures}. This proves both directions
globally, including all multiplicities and local unit coefficients.

For $k=1$, this is simply $h(\lambda)=0$ at every distinct zero,
without changing or discarding its original multiplicity in $\Xi$.
The derivative ambiguity has the exact presentation
\begin{equation}\label{eq:cq-firstambiguity}
 \cqE/\{h:h(\lambda)=0\text{ for every }\Xi(\lambda)=0\}
 \ \xrightarrow{\ \cong\ }\ (DI+I)/I,
 \qquad [h]\longmapsto[\Xi'h].
\end{equation}
Indeed $D(\Xi h)=\Xi h'+\Xi'h$, and the first term lies in $I$.
At each zero of multiplicity $m$, $\Xi'$ has exactly order $m-1$,
so its product with $h$ is divisible by $\Xi$ exactly when
$h(\lambda)=0$ for every zero. This proves the precise kernel,
image, and inverse on the image in \eqref{eq:cq-firstambiguity}.

For the heat generator ($k=2$), at any zero of multiplicity at least
two the condition is $h(\lambda)=h'(\lambda)=0$.
At a simple zero it is instead
\begin{equation}\label{eq:cq-simpleheatdomain}
 A_\lambda(\lambda)h'(\lambda)
        +A_\lambda'(\lambda)h(\lambda)=0.
\end{equation}
It does not require $h(\lambda)=0$. These follow from the index
range in \eqref{eq:cq-Ik}: the simple zero has only the coefficient
$b_1$ constrained. In particular the second-derivative relation
domain need not be a multiplication module: at a simple zero, a
local $\Xi h$ obeying \eqref{eq:cq-simpleheatdomain} and with
$h(\lambda)\ne0$ has
$D^2(s\Xi h)(\lambda)=2(\Xi h)'(\lambda)\ne0$.
This is a local statement with its explicit witness condition; it
does not assume the existence of a global simple zeta zero.

There is also an exact calculation of every finite projection of
the ambiguity $W_k=(D^kI+I)/I$. In the factorial jet coordinates
of any finite set $F$ of distinct zeros, its image is precisely
\begin{equation}\label{eq:cq-ambiguityjets}
 \prod_{\lambda\in F}
 \epsilon^{\max(m_\lambda-k,0)}
       \cqC[\epsilon]/(\epsilon^{m_\lambda}).
\end{equation}
The local dimension is $\min(k,m_\lambda)$. The differentiated
series above proves containment. For surjectivity, prescribe the
allowed output coefficient $y_j$ of $z^j$, where
$\max(m-k,0)\leq j<m$. Set
$b_{j+k-m}=j!y_j/(j+k)!$ and give the other coefficients
$b_0,\ldots,b_{k-1}$ arbitrary values, say zero. Finite Taylor
division by the nonzero unit $A_\lambda$ then specifies the first
$k$ Taylor coefficients of $h$ at $\lambda$.
These finitely many specifications can be achieved by one polynomial:
put $M_\lambda(s)=\prod_{\nu\in F,\nu\ne\lambda}(s-\nu)^k$;
take the Taylor polynomial of degree $k-1$ of the prescribed local
polynomial divided by $M_\lambda$, and sum the resulting
$M_\lambda$ multiples over $F$. At its own node each summand
agrees with the prescribed polynomial modulo $(s-\lambda)^k$,
and at the other nodes it vanishes to that order. This constructs
the required $h\in\cqC[s]\subset\cqE$ and proves
\eqref{eq:cq-ambiguityjets}. No claim about arbitrary infinite jet
interpolation has replaced this exact finite image.

\begin{cqtheorem}[Direct differential and heat obstructions]
\label{thm:cq-obstruction}
Neither $N_\gamma$ nor $N_\kappa$ is invariant under $D$ or $D^2$.
For either frozen quotient, neither representative formula $[f']$
nor $[-f''/4]$ defines an operator on any nonempty class domain
admitting all representatives.
For every $t\ne0$, neither $U_tN_\gamma\subseteq N_\gamma$ nor
$U_tN_\kappa\subseteq N_\kappa$ holds. Ordinary multiplication
by $s$ preserves neither $U_tN_\gamma$ nor $U_tN_\kappa$.
\end{cqtheorem}
\begin{proof}
Choose any zero $\lambda$ of the actual $\Xi$ and let $m\geq1$
be its multiplicity. The existence of its zeros is part of the
classical arithmetic datum in the primary source; no realness or
simplicity is used. Write $\Xi=(s-\lambda)^mA$ with
$A(\lambda)\ne0$. Then $\Xi'=(s-\lambda)^{m-1}
(mA+(s-\lambda)A')$ does not vanish to order $m$.
Since $\Xi\in N_\gamma\subseteq N_\kappa$, both derivative
invariance assertions fail. If $m\geq2$, $\Xi''$ has leading
coefficient $m(m-1)A(\lambda)$ at order $m-2$, proving the
second-derivative failure. If $m=1$, use the actual polynomial
multiple $n=(s-\lambda)\Xi$ instead: $n''(\lambda)=2\Xi'(\lambda)
\ne0$. This retains the product-rule cross term and covers all
multiplicities. The nonempty-domain assertion follows from the
explicit fibre computation in \eqref{eq:cq-relation}.

For $a\in\cqC$, $e^{as}\Xi\in N_\gamma$ by polynomial
approximation of $e^{as}$ and continuous multiplication, and hence
belongs to $N_\kappa$. The conjugation $D(e^{as}f)=e^{as}(D+a)f$
and the absolutely convergent heat and translation series give
\begin{equation}\label{eq:cq-exponential}
 U_t(e^{as}\Xi)(s)
    =e^{as-ta^2/4}Z_t(s-ta/2).
\end{equation}
For completeness, the translation series for $f$ is its Taylor
series about $s$ and sums to $f(s-ta/2)$; its local uniform
convergence and the derivative estimates justify commutation of
$e^{-(ta/2)D}$ with $U_t$. If $U_tN_\bullet\subseteq N_\bullet$
for either choice, evaluate \eqref{eq:cq-exponential} at $\lambda$.
The exponential never vanishes, so $Z_t(\lambda-ta/2)=0$ for
every $a$. When $t\ne0$ these arguments fill $\cqC$; this contradicts
$Z_t\ne0$. Finally if $M$ preserved $U_tN_\bullet$, conjugation
would imply $(M+(t/2)D)N_\bullet\subseteq N_\bullet$.
Since $M$ already preserves $N_\bullet$, this forces
$DN_\bullet\subseteq N_\bullet$, contrary to the first part.
\end{proof}

The analytic moving-divisor quotient, which does carry the zeros of
the actual $Z_t$ under ordinary $s$, is
\begin{equation}\label{eq:cq-divisor}
 I_t=\overline{\cqC[s]Z_t}^{\,\kappa}
 =\{f\in\cqE:f/Z_t\text{ is entire}\},\qquad
 A_t=(\cqE,\kappa)/I_t.
\end{equation}
The division lemma and polynomial approximation further prove
\[
 I_t=Z_t\cqE=\overline{\cqC[s]Z_t}^{\,\gamma}\quad(t\in\cqC).
\]
For a complex $t$ with $Z_t(0)=0$, use the retained nonzero center
$c$ in the division lemma; $Z_t$ is not identically zero.
The analytically transported relation has the different exact
principal presentation for the proved transported multiplication:
\begin{equation}\label{eq:cq-starideal}
 U_t(\Xi\cqE)=Z_t\star_t\cqE.
\end{equation}
Indeed $U_t(\Xi h)=Z_t\star_t U_th$, and $U_t$ is onto.
The map $h\mapsto Z_t\star_t h$ is a continuous linear
isomorphism of $(\cqE,\gamma)$ onto this closed relation;
its inverse on that relation is
$n\mapsto U_t((U_{-t}n)/\Xi)$, continuous by the division and
heat lemmas. Thus the two multiplications and their exact relation
spaces are connected by a proved isomorphism with a proved inverse.
The proof of Theorem~\ref{thm:cq-compact} applies verbatim with
the nonzero entire order-at-most-one function $Z_t$: it uses only
local Taylor division and compact polynomial approximation.
Consequently $\cqSpec(M\text{ on }A_t)=\{\lambda:Z_t(\lambda)=0\}$,
with the full local multiplicity jets, and with the explicit inverse
$[f]\mapsto[(f-f(\lambda)Z_t/Z_t(\lambda))/(s-\lambda)]$
off that set. Equation \eqref{eq:cq-zero} is the evolution of this
original-coordinate spectrum on its simple-zero locus.
This spectrum can move; \eqref{eq:cq-specmove} is the spectrum
of the different, exactly specified transported operator $B_t$.

We now give the promised comparison morphisms with every kernel
retained. Put $J=\cqN\cap N_\kappa$ and
$C=(\cqE,\tau\vee\kappa)/J$, where $\tau\vee\kappa$ is the
topology generated by both topologies. Since each summand is closed
in its own topology, $J$ is closed in their join. The identity maps
give continuous surjections
\begin{equation}\label{eq:cq-span}
 \cqQ\ \longleftarrow\ C\ \longrightarrow\ Q_\kappa,
 \qquad\ker(C\to\cqQ)=\cqN/J,
 \quad\ker(C\to Q_\kappa)=N_\kappa/J.
\end{equation}
Each assertion follows by taking the same representative $f$ and
testing membership in the respective relation subspace; the quotient
universal property proves continuity. An algebraic common quotient is
$\cqE/(\cqN+N_\kappa)$ with the quotient topology from the join;
its Hausdorff version uses
$\overline{\cqN+N_\kappa}^{\,\tau\vee\kappa}$.
The map from $C$ to this further quotient is continuous. Continuity
of maps to it from each of the two separately topologized quotients
is not asserted. No closedness of the sum is presumed. This describes
the precise information lost in that further quotient.

At time $t$ replace the two subspaces by $N_t$ and $I_t$ and the
join by $\tau_t\vee\kappa$. The same construction gives continuous
surjections from their intersection quotient to the genuine transported
arithmetic quotient $Q_t$ and to $A_t$, with kernels
$N_t/(N_t\cap I_t)$ and $I_t/(N_t\cap I_t)$, respectively.
The two generator systems are linked by the explicit formula
$R(B_t)Z_t$ versus $R(s)Z_t$ and the defects
\eqref{eq:cq-generators}--\eqref{eq:cq-unit}.
These are exact morphisms and closure formulas for the actual source
space. Determining that either kernel vanishes for $\tau$ is not
established by the two inspected primary texts. Accordingly the
source-frozen statements $D\cqN\subseteq\cqN$ and
$D^2\cqN\subseteq\cqN$ have not been decided solely from that
unnamed topology. The unqualified noninvariance results above have
been proved in the two explicitly specified analytic realizations.

\section{Arithmetic traces and their exact Mellin obstructions}
\label{sec:cq-arithmetic}
\subsection{The full Schwartz trace and arithmetic principal parts}
\label{sec:cq-principalparts}
Retain the source domain
\[
 \mathcal S_0^{\rm ev}
 =\{\psi\in\mathcal S(\mathbb R):
       \psi(-x)=\psi(x),\ \psi(0)=\widehat\psi(0)=0\},
 \qquad
 \widehat\psi(y)=\int_{\mathbb R}\psi(x)e^{-2\pi ixy}\,dx.
\]
Its full trace map takes values in entire functions:
\begin{equation}\label{eq:cq-Schwartz-trace}
 \begin{aligned}
 F_\psi(s)&=\mathcal F_\mu\mathcal E\psi(s),\\
 F_\psi(i(z-\tfrac12))&=\zeta(z)\mathcal M\psi(z),
 &\mathcal M\psi(z)&=\int_0^\infty\psi(x)x^z\,\frac{dx}{x}.
 \end{aligned}
\end{equation}
The Mellin integral is holomorphic for $\operatorname{Re}z>-2$,
and the factorization holds there with its removable value at
$z=1$. This statement does not impose an unproved order bound
on the entire $F_\psi$ for arbitrary Schwartz inputs.

Here is a direct proof with all analytic domains retained.
For fixed $x>0$, periodize $\psi$ by the smooth function
\[
 y\longmapsto\sum_{k\in\mathbb Z}\psi(x(k+y)).
\]
Its Fourier coefficients are $x^{-1}\widehat\psi(\ell/x)$:
integrate on $y\in[0,1]$, substitute $v=x(k+y)$, and sum the
adjacent intervals, justified by Schwartz decay. These coefficients
decrease faster than every power of $\ell$, so their Fourier series
converges with all derivatives and at $y=0$ gives the Poisson
identity. Evenness and the two exact zero conditions yield
\[
 \sum_{n\geq1}\psi(nx)=x^{-1}\sum_{n\geq1}\widehat\psi(n/x),
 \qquad \mathcal E\psi(x)=\mathcal E\widehat\psi(1/x).
\]
For $x\geq1$, a Schwartz bound of arbitrarily high order gives
$|\sum\psi(nx)|\leq C_Ax^{-A}\sum n^{-A}$ for $A>1$.
The same estimate after termwise differentiation follows by choosing
the Schwartz exponent larger than the derivative order plus one.
The Poisson identity supplies the corresponding bounds at zero.
Thus $\mathcal E\psi(e^u)$ decreases faster than $e^{-A|u|}$
for every $A>0$ at both ends. Its Fourier integral and each
derivative in $s$ are dominated locally uniformly in $s$ by an
integrable function. This proves that $F_\psi$ is entire.

Since $\psi$ is even and $\psi(0)=0$, its Taylor formula gives
$\psi(x)=O(x^2)$ at zero. Together with Schwartz decay at infinity
this proves the stated holomorphy of $\mathcal M\psi$ by
dominated differentiation, including all powers of $\log x$.
For $\sigma=\operatorname{Re}z>1$,
\[
 \sum_{n\geq1}\int_0^\infty|\psi(nx)|x^\sigma\,\frac{dx}{x}
 =\Bigl(\sum_{n\geq1}n^{-\sigma}\Bigr)
          \int_0^\infty|\psi(y)|y^\sigma\,\frac{dy}{y}<\infty.
\]
Absolute interchange and the original exponent $x^{-is}$ now give
\eqref{eq:cq-Schwartz-trace} on this half-plane. Moreover
$\mathcal M\psi(1)=\int_0^\infty\psi(x)\,dx
=\widehat\psi(0)/2=0$. This cancels the simple pole of $\zeta$
at $1$, so the right side is holomorphic for
$\operatorname{Re}z>-2$. The identity theorem extends the equality
to that entire half-plane, proving the assertion.

Let $\mathcal Z$ denote the set of distinct nontrivial zeros $\rho$
of $\zeta$, with their original multiplicities $m_\rho$.
The evenness of $\Xi$, immediate from the retained cosine integral,
gives the exact reflected-coordinate identity
\begin{equation}\label{eq:cq-reflected-Xi}
 s=i(z-\tfrac12),\qquad z=\tfrac12-is,\qquad
 \Xi(i(z-\tfrac12))=\xi(1-z)=\xi(z).
\end{equation}
At each $\rho\in\mathcal Z$ the completion factor
$z(z-1)\pi^{-z/2}\Gamma(z/2)/2$ is a nonvanishing holomorphic
unit. Hence $\lambda_\rho=i(\rho-\tfrac12)$ is a zero of
$\Xi$ of exactly multiplicity $m_\rho$, and these are all its
zeros. No reflection of the retained spectral operator $s$ has
been performed: \eqref{eq:cq-reflected-Xi} is the exact Mellin
parameter dictionary for that same operator. The factorization
\eqref{eq:cq-Schwartz-trace} proves
\begin{equation}\label{eq:cq-fulltrace-jets}
 \begin{gathered}
 D^jF_\psi(\lambda_\rho)=0
 \quad(\rho\in\mathcal Z,\ 0\leq j<m_\rho),\\
 F_\psi\in\cqE\ \Longrightarrow\ F_\psi\in\Xi\cqE.
 \end{gathered}
\end{equation}
The first assertion follows because $\mathcal M\psi$ is holomorphic
at every $\rho$ and $s-\lambda_\rho=i(z-\rho)$; the second
then follows from the proved division theorem. It is an exact
intersection assertion about the full Schwartz trace, without
enlarging the order-at-most-one domain or closing it in $\tau$.

For each $\rho$ retain the actual local unit
$a_\rho(z)=\zeta(z)/(z-\rho)^{m_\rho}$ and the space
\[
 \mathcal P_\rho^-=
  \left\{\sum_{j=1}^{m_\rho}c_{-j}(z-\rho)^{-j}:c_{-j}\in\cqC\right\}.
\]
The principal-part map on the original entire-function set is
\begin{equation}\label{eq:cq-principalparts}
 \mathfrak P:\cqE\longrightarrow\prod_{\rho\in\mathcal Z}\mathcal P_\rho^-,
 \qquad
 F\longmapsto\left(\operatorname{PP}_{z=\rho}
                 \frac{F(i(z-\tfrac12))}{\zeta(z)}\right)_\rho.
\end{equation}
The kernel of this map is exactly $\Xi\cqE$. It induces an
injective linear map from the underlying vector space of each
$Q_\gamma,Q_\kappa$ onto the same specified image. Every finite
projection is surjective. The map from $Q_\kappa$ to its image,
with the product subspace topology, is a topological linear
isomorphism; the product of all principal-part spaces is its
completion. From $Q_\gamma$ it is a continuous bijection onto
that image whose inverse is not continuous.

To prove each claim while retaining the analytic units, write
$a_\rho(\rho+w)^{-1}=\sum_{\ell\geq0}\alpha_{\rho,\ell}w^\ell$.
The exact Laurent coefficient at degree $-m_\rho+q$, $0\leq q<m_\rho$,
of the $\rho$ component is
\begin{equation}\label{eq:cq-principalpart-coefficients}
 c_{-m_\rho+q}
 =\sum_{j=0}^{q}\frac{i^j}{j!}F^{(j)}(\lambda_\rho)
                                  \alpha_{\rho,q-j}.
\end{equation}
This follows by multiplying the two convergent local Taylor
series and the retained factor $w^{-m_\rho}$. Its finite matrix
is triangular with nonzero diagonal entries
$i^j/(j!a_\rho(\rho))$. Explicitly, given a prescribed principal
part $P_\rho(w)$, take the coefficients $b_j$ through degree
$m_\rho-1$ of $a_\rho(\rho+w)w^{m_\rho}P_\rho(w)$;
the required original derivatives are $j!i^{-j}b_j$.
Multiplying back by $a_\rho^{-1}$ recovers the prescribed negative
powers, since the discarded error is divisible by $w^{m_\rho}$.
This proves a finite-dimensional linear isomorphism at every zero,
including its inverse and every factorial and power of $i$.

The vanishing of all principal parts is therefore equivalent to
the vanishing of all the original multiplicity jets of $F$.
Theorem~\ref{thm:cq-equalclosures} identifies that kernel exactly
with $\Xi\cqE$. At any finite set of zeros, the required derivatives
can be realized by the explicit Hermite lift
\eqref{eq:cq-finite-Hermite-lift}; this proves surjectivity of
every finite principal-part projection. Each block map and its
inverse are continuous because their finite matrices have fixed
complex coefficients. Their product is therefore a homeomorphism
of the two product spaces: the inverse image of a cylinder fixing
finitely many coordinates is again controlled by those finitely
many blocks. Applying this product isomorphism to
Theorem~\ref{thm:cq-jet-topology} proves all stated topology,
completion, and comparison assertions for $\mathfrak P$.

This map also carries the original spectral action exactly. For
$P_\rho=\mathfrak P_\rho(F)$ one has
\begin{equation}\label{eq:cq-s-principalparts}
 \mathfrak P_\rho(sF)
 =\operatorname{PP}_{w=0}
                  ((\lambda_\rho+iw)P_\rho(w)).
\end{equation}
Indeed the full Laurent quotient for $sF$ is the quotient for $F$
multiplied by $i(z-\tfrac12)=\lambda_\rho+iw$.
Its holomorphic residual contributes no negative power after this
multiplication. Thus the finite block is the retained
$\lambda_\rho$ plus the explicitly specified nilpotent operator
$P\mapsto\operatorname{PP}(iwP)$, whose nilpotency index is
$m_\rho$: iteration on $w^{-m_\rho}$ gives a nonzero multiple
of $w^{-1}$ at step $m_\rho-1$ and zero at step $m_\rho$.

For $I=\Xi\cqE$, the derivative ambiguity
$W_k=(D^kI+I)/I$ has, under $\mathfrak P$, the exact finite image
\begin{equation}\label{eq:cq-principal-ambiguity}
 \prod_{\rho\in F}
 \left\{\sum_{j=1}^{\min(k,m_\rho)}c_{-j}(z-\rho)^{-j}:c_{-j}\in\cqC\right\}
 \quad(F\subset\mathcal Z\text{ finite}).
\end{equation}
The finite image in \eqref{eq:cq-ambiguityjets} consists exactly
of numerator jets divisible by $w^{\max(m_\rho-k,0)}$ after
the invertible substitution $s-\lambda_\rho=iw$.
Multiplication by the nonvanishing unit $a_\rho^{-1}$ preserves
this divisibility and is invertible on the truncated Taylor ring.
Division by $w^{m_\rho}$ therefore yields exactly all negative
powers of order at most $\min(k,m_\rho)$, proving both
containment and surjectivity in the displayed finite image.
In particular the complete local-unit expression, rather than only
a count of vanishing orders, transports the derivative defect to
the arithmetic Mellin poles.

\subsection{An exact arithmetic inverse for every heated Hermite trace}
\label{sec:cq-moebius}
Let $\psi$ be any finite linear combination of the original
$\psi_\ell^\pm$, retaining all its coefficients, and write
$F_0=\mathcal F_\mu\mathcal E\psi=p\Xi$. For $x>0$ and
$t\in\cqC$ define
\begin{equation}\label{eq:cq-moebius-data}
 \begin{aligned}
 q(x)&=\sum_{n\geq1}\psi(nx),&
 g_t(x)&=e^{t(\log x)^2/4}\mathcal E\psi(x),\\
 q_t(x)&=x^{-1/2}g_t(x)=e^{t(\log x)^2/4}q(x),&
 \psi_t(x)&=\sum_{n\geq1}\mu(n)q_t(nx).
 \end{aligned}
\end{equation}
Here $\mu(1)=1$, $\mu(n)=0$ when a prime square divides $n$,
and $\mu(n)=(-1)^j$ for a product of $j$ distinct primes.
For the inverse $\psi_t$, the symbol $\mathcal E\psi_t$ means
the absolutely convergent trace-sum extension
$x^{1/2}\sum_{m\geq1}\psi_t(mx)$ on the positive half-line.
Its domain here consists of smooth functions with rapid decrease
at infinity; no Schwartz extension through zero is required.
The transform and inverse assertions, with their domains, are
\begin{equation}\label{eq:cq-moebius-inverse}
 \mathcal E\psi_t=g_t,\qquad
 \mathcal F_\mu g_t=U_t(p\Xi)=:F_t,\qquad
 \int_0^\infty\psi_t(x)x^z\,\frac{dx}{x}
       =\frac{F_t(i(z-\tfrac12))}{\zeta(z)}
       \quad(\operatorname{Re}z>1).
\end{equation}
The series for $\psi_t$ converges locally with all derivatives on
$(0,\infty)$, and $\psi_t$ and all its derivatives decrease faster
than every inverse power at infinity. No behavior at zero has been
assumed for this inverse.

We prove all analytic assertions in these formulas. A finite Hermite
combination and its Fourier transform are polynomials times the
original Gaussian $e^{-\pi x^2}$. Thus at infinity, for every
derivative order $j$, their derivatives are bounded by
$C_j(1+x)^{d_j}e^{-\pi x^2}$. For $x\geq1$, termwise derivatives
of $\sum\psi(nx)$ are bounded by a sum of
$C_j n^j(1+nx)^{d_j}e^{-\pi n^2x^2}$; splitting the Gaussian
into two equal exponential factors bounds that sum by
$C'_j(1+x)^{d_j}e^{-\pi x^2/2}$. This also justifies the
differentiation. Both zero conditions give, by the Poisson identity,
\[
 \mathcal E\psi(x)=\mathcal E\widehat\psi(1/x).
\]
Consequently, in $u=\log x$, the functions
$\mathcal E\psi(e^u)$ and all their $u$ derivatives have bounds
of the form $C\exp(C'|u|-c e^{2|u|})$ on the two ends.
Multiplication by $e^{tu^2/4}$, and any finite number of its
derivatives, preserves integrability against $e^{A|u|}$ for every
$A>0$, locally uniformly for $t$ in compact subsets of $\cqC$.
It follows that $\mathcal F_\mu g_t$ is entire in both $t,s$.
Since $D_s^{2k}e^{-isu}=(-1)^ku^{2k}e^{-isu}$, absolute
termwise integration of the time exponential gives exactly
\[
 \int_{\mathbb R}e^{tu^2/4}\mathcal E\psi(e^u)e^{-isu}\,du
 =\sum_{k\geq0}\frac{(-t/4)^k}{k!}
       D_s^{2k}(p(s)\Xi(s))=F_t(s).
\]
The heat lemma identifies the same series in $\gamma$; hence
$F_t\in\cqE$ with the original coordinate and coefficient.

The proved estimates imply that $q_t$ and all its $x$ derivatives
decrease faster than every power at infinity, and that $q_t$
decreases faster than every power at zero. On a compact interval
$x\in[\delta,L]\subset(0,\infty)$, choose $A>j+1$ to bound
the $j$th derivative of the $n$th inverse summand by
$C n^j(nx)^{-A}\leq C\delta^{-A}n^{j-A}$.
The finitely many arguments below one are harmless. This proves
local convergence with derivatives. On $x\geq1$ the same bound,
with arbitrarily large $A$, proves rapid decrease of each derivative.
For a fixed $x>0$ the double sum used below is absolutely bounded by
\[
 \sum_{m,n\geq1}|q_t(mnx)|
 =\sum_{k\geq1}d(k)|q_t(kx)|<\infty,
\]
where $d(k)\leq k$ and rapid decrease permits an exponent greater
than two. Unique prime factorization gives
$\sum_{n\mid k}\mu(n)=1$ for $k=1$ and
$\prod_{p\mid k}(1-1)=0$ for $k>1$. Therefore
\[
 \sum_{m\geq1}\psi_t(mx)
 =\sum_{k\geq1}\Bigl(\sum_{n\mid k}\mu(n)\Bigr)q_t(kx)
 =q_t(x).
\]
Multiplication by the retained factor $x^{1/2}$ proves the first
identity of \eqref{eq:cq-moebius-inverse}. It also proves uniqueness
of this inverse among functions decreasing faster than every power
at infinity: for any such $v$ with $\sum_m v(mx)=q_t(x)$, the
same absolutely convergent double sum applied to
$\sum_n\mu(n)\sum_m v(mnx)$ equals $v(x)$.

For $\sigma=\operatorname{Re}z>1$, absolute integration gives
\[
 \sum_{n\geq1}|\mu(n)|\int_0^\infty
       |q_t(nx)|x^\sigma\,\frac{dx}{x}
 \leq\Bigl(\sum_{n\geq1}n^{-\sigma}\Bigr)
       \int_0^\infty|q_t(y)|y^\sigma\,\frac{dy}{y}<\infty.
\]
The divisor identity also gives
$(\sum\mu(n)n^{-z})\zeta(z)=1$ by absolute multiplication of
the two Dirichlet series in this half-plane. Finally
\[
 \int_0^\infty q_t(x)x^z\,\frac{dx}{x}
 =\int_0^\infty g_t(x)x^{z-1/2}\,\frac{dx}{x}
 =F_t(i(z-\tfrac12)),
\]
because $-i\,i(z-\tfrac12)=z-\tfrac12$.
These equalities prove the last identity of
\eqref{eq:cq-moebius-inverse}, including its absolute domain.

At a nontrivial zero $\rho$ of $\zeta$ of multiplicity $m$, put
$\lambda=i(\rho-\tfrac12)$ and retain
$\zeta(z)=(z-\rho)^m a_\rho(z)$ with $a_\rho(\rho)\ne0$.
The retained cosine integral makes $\Xi$ even, so its definition gives
\[
 \Xi(i(z-\tfrac12))=\xi(1-z)=\xi(z).
\]
At a nontrivial zero, the factor
$z(z-1)\pi^{-z/2}\Gamma(z/2)/2$ is a nonvanishing holomorphic
unit. Thus this reflected coordinate sends the zero at $\rho$
to the zero at $\lambda$ with exactly the same multiplicity.
The meromorphic continuation of the last quotient has pole order
\begin{equation}\label{eq:cq-moebius-pole}
 \max\{m-\operatorname{ord}_\lambda F_t,0\}.
\end{equation}
For the zero seed $F_t=0$, the quotient is zero and its pole order
is zero; the displayed formula uses $\operatorname{ord}_\lambda0=+\infty$.
Indeed $s-\lambda=i(z-\rho)$, so its numerator has exactly
the original order $\operatorname{ord}_\lambda F_t$, with its
leading coefficient multiplied by the nonzero power of $i$.
The unit $a_\rho$ is retained in every principal coefficient.

For every fixed $t\ne0$ and every zero $\lambda$ of $\Xi$, there
is a polynomial $p$ for which $F_t(\lambda)\ne0$: otherwise all
$U_t(s^n\Xi)$ would vanish there, contrary to the no-common-zero
proof in Section~\ref{sec:cq-moving}. The exact polynomial-image
calculation supplies a finite original Hermite seed for this $p$.
For that seed the unique inverse $\psi_t$ in
\eqref{eq:cq-moebius-data} is unbounded near zero. To prove this
last assertion, boundedness near zero together with its proved rapid
decrease at infinity would make its Mellin integral holomorphic in
$\operatorname{Re}z>0$: on compact subsets, derivatives in $z$
are dominated by powers of $|\log x|x^{\sigma-1}$ near zero
with $\sigma>0$. Its equality in $\operatorname{Re}z>1$ would
then continue as a holomorphic function through every nontrivial
$\rho$, contradicting the pole of order $m$ just proved at
$\rho=\tfrac12-i\lambda$. Thus for each nonzero heat time
some actual finite Hermite trace has an explicit arithmetic inverse
which fails even boundedness at the origin. This proves a concrete
obstruction to Schwartz descent with a unique inverse and its poles;
it does not replace the unresolved closure in $\tau$ by a different
closure.


\section{The inverse-root cover, the branch locus, and arithmetic fibres}
\label{sec:cq-cover}
The retained polynomial inverse-fibre slice is
$P_{a,0,0}(r)=-2r^2-2a$. Its root equation is $a=-r^2$.
The arithmetic coordinate is $a=2s$, so
\begin{equation}\label{eq:cq-cover}
 s=-r^2/2.
\end{equation}
The finite free extension of the actual quotient is
\[
 \widehat Q=\cqC[r]\otimes_{\cqC[s]}\cqQ
           =\cqQ\oplus r\cqQ,
 \qquad
 \mathcal R=\begin{pmatrix}0&-2S\\1&0\end{pmatrix},
 \qquad -\tfrac12\mathcal R^2=\begin{pmatrix}S&0\\0&S\end{pmatrix}.
\]
The even--odd decomposition of a polynomial in $r$ proves freeness
with basis $(1,r)$ and uniqueness coefficient by coefficient.
Multiplying $q_0+rq_1$ by $r$ proves the displayed matrices;
the finite direct-sum topology makes every induced map continuous.
The involution is $(q_0,q_1)\mapsto(q_0,-q_1)$; the trace is
$\operatorname{Tr}(q_0+rq_1)=2q_0$, with kernel $r\cqQ$.
It is $\cqC[s]$-linear. It does not provide a $\cqC[r]$-module
quotient with this kernel unless $S=0$: multiplying $rq$ by $r$
gives $-2Sq$ in the even summand. This proves the exact limitation,
rather than suppressing the odd component. The same construction
at time $t$ uses $S_t$ and is conjugated by
$\operatorname{diag}(\mathcal U_t,\mathcal U_t)$.

The original $s$-spectrum is unchanged by the extension: an inverse
of $S-\lambda$ duplicates to a diagonal inverse, and conversely
projection of an inverse of the doubled matrix to its first component,
following inclusion of that component, gives an inverse of
$S-\lambda$. Both directions are continuous. Moreover
\begin{equation}\label{eq:cq-rspec}
 \cqSpec(\mathcal R)
 =\{\rho\in\cqC:-\rho^2/2\in\cqSpec(S)\}.
\end{equation}
To prove this without an unsupported spectral-mapping citation,
solve $(\mathcal R-\rho)(u,v)=(a,b)$. The second row gives
$u=b+\rho v$ and the first becomes
$-2(S+\rho^2/2)v=a+\rho b$. A continuous inverse of
$S+\rho^2/2$ therefore gives a continuous inverse matrix.
Conversely the right-hand side $(a,0)$ and the second projection
of an inverse matrix supply an inverse of $-2(S+\rho^2/2)$,
with both identities checked by the same two rows.

Before any quotient by $\Xi$, let $\mathcal O_s$ be the sheaf of
holomorphic functions in $s$ and put
$B=\mathcal O_s[r]/(r^2+2s)$, identified with
$\mathcal O_s\oplus r\mathcal O_s$. On the punctured base $s\ne0$,
implicit differentiation gives
$dr/ds=-1/r=r/(2s)$. Thus the actual lifted derivative is
\begin{equation}\label{eq:cq-connection}
 \nabla_s\binom{f_0}{f_1}
 =\binom{f_0'}{f_1'+f_1/(2s)},\quad
 R=\begin{pmatrix}0&-2s\\1&0\end{pmatrix},\quad
 [\nabla_s,R]=-R^{-1}.
\end{equation}
Indeed with $A=\operatorname{diag}(0,1/(2s))$,
$R'+AR-RA=\left(\begin{smallmatrix}0&-1\\1/(2s)&0\end{smallmatrix}\right)
=-R^{-1}$. A second application gives
\begin{equation}\label{eq:cq-connection2}
 \nabla_s^2=\operatorname{diag}
 \left(D^2,D^2+s^{-1}D-\frac1{4s^2}\right).
\end{equation}
Equivalently, on a function of $r$ the derivative is
$-r^{-1}\partial_r$ and the heat generator is
$-\tfrac14r^{-2}\partial_r^2+\tfrac14r^{-3}\partial_r$.
Squaring $-r^{-1}\partial_r$ by the product rule proves this
with both signs preserved.

The branch extension is exact. The meromorphic connection sends
$B$ into $s^{-1}B$ and the connection one-form has a simple pole
with residue $\operatorname{diag}(0,1/2)$. Its maximal domain for
holomorphic output at $s=0$ is
\begin{equation}\label{eq:cq-branchdomain}
 \operatorname{Dom}_{\rm hol}\nabla_s
       =\mathcal O_s\oplus r s\mathcal O_s.
\end{equation}
In fact the odd coefficient is $f_1'+f_1/(2s)$; its only possible
pole has coefficient $f_1(0)/2$, proving necessity and sufficiency.
For the square, write $f_1=\sum_{n\geq0}a_ns^n$. Its image has
coefficient $(n^2-1/4)a_ns^{n-2}$. The negative powers vanish
exactly when $a_0=a_1=0$, so
\[
 \operatorname{Dom}_{\rm hol}\nabla_s^2
       =\mathcal O_s\oplus r s^2\mathcal O_s.
\]
For $k$ iterates the odd coefficient is
$\prod_{h=0}^{k-1}(n+1/2-h)a_ns^{n-k}$, whose prefactor is
never zero for integral $n$. Thus its domain is
$\mathcal O_s\oplus r s^k\mathcal O_s$.
No derivation $B\to B$ extending $\partial_s$ exists at the branch:
derivation of $r^2+2s=0$ would give $2r\delta(r)+2=0$,
forcing $r$ to be a unit although it vanishes at the branch point.
The logarithmic derivative $s\nabla_s$
does extend to all of $B$, with matrix
$sD+\operatorname{diag}(0,1/2)$.

These are sheaf and representative calculations; they do not assert
that $\nabla_s$ descends through the frozen arithmetic relation.
For a holomorphic local multiplier $R_0(s)$ the relation vector
$(R_0\Xi,0)$ acquires $(R_0'\Xi+R_0\Xi',0)$, and the odd
relation $(0,R_0\Xi)$ acquires
$(0,R_0'\Xi+R_0\Xi'+R_0\Xi/(2s))$.
The second iterate includes the full $2R_0'\Xi'$ term, as well as
the odd $s^{-1}D-1/(4s^2)$ terms of
\eqref{eq:cq-connection2}. Jet extension
\eqref{eq:cq-jetextension}, or the moving closed quotient
\eqref{eq:cq-Qt}, retains these relations exactly.

Finally, the branch point is not an arithmetic spectral point at
$t=0$. The source gives $\Xi(0)\ne0$, so its Proposition 3.6
makes $S$ continuously invertible; then $\mathcal R$ is invertible
by \eqref{eq:cq-rspec}. In the analytic divisor realization the
stronger local statement follows from $\Xi$ being a holomorphic
unit on a disk about zero: the quotient stalk is zero there.
For real $t$, $Z_t(0)>0$ directly from the retained kernel:
\[
 2\pi^2n^4e^{9u}-3\pi n^2e^{5u}
 =\pi n^2e^{5u}(2\pi n^2e^{4u}-3)>0\quad(u\geq0),
\]
and the positive integral converges by the bound already proved.
Thus the actual real heat divisor never meets the cover branch
point $s=0$. On a simple zero curve its retained lift satisfies
$r^2=-2\lambda$, $r'=-\lambda'/r$; together with
\eqref{eq:cq-zero} these are the original-coordinate and cover
equations, including the unit contribution. A nonreal $r$ alone
does not specify whether $\lambda=-r^2/2$ is nonreal.

\section{Exact ingress from compact Weil tests}
\label{sec:cq-tests}
The accompanying retained Weil-test calculation supplies the concrete
space $\mathcal E=C_c^\infty(\mathbb R;\cqC)$ with coordinate
operator $T=-i\partial_u$ and Fourier transform
\[
 (\mathscr Ff)(s)=\int_{\mathbb R}f(u)e^{-ius}\,du.
\]
Here we give the complete heat ingress and its inverse domains;
none of the following claims requires identifying the source
quotient topology with a test-space topology. Integration by parts
gives $\mathscr F(Tf)=s\mathscr Ff$ with no endpoint term.
If $\operatorname{supp}f\subset[-L,L]$, then
\[
 |\mathscr Ff(s)|\leq\|f\|_1e^{L|\operatorname{Im}s|},
 \qquad
 \|\mathscr Ff\|_{p,a}\leq
 \|f\|_1\sup_{r\geq0}e^{Lr-ar^p}<\infty.
\]
Thus the Fourier transform maps the actual test space into the
original entire-function set, continuously in $\gamma$ on every
fixed-support smooth-function space and hence in its usual
inductive-limit topology.

Define, for all complex $t$,
\begin{equation}\label{eq:cq-testheat}
 (G_tf)(u)=e^{tu^2/4}f(u),\qquad
 U_t\mathscr Ff=\mathscr F G_tf,\qquad
 T_t=G_tTG_{-t}=T+\frac{it}{2}u.
\end{equation}
The multiplication $G_t$ is a continuous automorphism on each
fixed-support smooth-function space, with inverse $G_{-t}$.
Indeed each derivative of the smooth multiplier is bounded on the
fixed compact support, and Leibniz's finite sum bounds every output
derivative seminorm by finitely many input ones. The same argument
for the inverse proves the assertion. The exponential series and
every $u$-derivative converge uniformly on a fixed compact set,
locally uniformly in $t$; consequently the family is entire in time
in that test topology. In the Fourier integral,
$D^{2k}\mathscr Ff=\mathscr F((-1)^ku^{2k}f)$.
The compact support permits absolute interchange with the
exponential series, proving the middle identity with the exact
factor $1/4$. Finally differentiation of $e^{-tu^2/4}f$ gives
the displayed $T_t$, and
\[
 \mathscr F(T_tf)=B_t\mathscr Ff,
 \qquad D\mathscr Ff=\mathscr F(-iu f).
\]
Both equations follow directly from the integral, including their
signs. The test-space version of the retained cover is therefore
\[
 \mathcal R_t^{\rm test}
   =\begin{pmatrix}0&-2T_t\\1&0\end{pmatrix},\qquad
 (\mathcal R_t^{\rm test})^2=-2\operatorname{diag}(T_t,T_t).
\]
Componentwise Fourier transformation intertwines this matrix with
the representative matrix using $B_t$, and then with the actual
quotient matrix using $S_t$.

To retain every possible arithmetic relation in this ingress, put
\begin{equation}\label{eq:cq-testkernel}
 K=\{f\in\mathcal E:\mathscr Ff\in\cqN\},\qquad
 K_t=\{f\in\mathcal E:\mathscr Ff\in N_t\}.
\end{equation}
Then $K_t=G_tK$: by \eqref{eq:cq-testheat},
$\mathscr Ff\in U_t\cqN$ exactly when
$\mathscr F(G_{-t}f)\in\cqN$. Thus there are injective linear
maps onto their explicitly specified images
\[
 \mathcal E/K\longrightarrow\cqQ,\quad[f]\longmapsto[\mathscr Ff],
 \qquad
 \mathcal E/K_t\longrightarrow Q_t,\quad[f]\longmapsto[\mathscr Ff],
\]
and the square formed by $G_t$ and $\mathcal U_t$ commutes.
The kernel statements follow exactly from
\eqref{eq:cq-testkernel}; no claim that $K=0$ in the unnamed
source topology is inserted. To make the ingress continuous without
that identification, give $\mathcal E$ the join $\eta$ of its
usual test topology with the inverse-image topology of $\tau$
under $\mathscr F$. Then $K$ is closed and the first quotient map
is continuous. Give the time-$t$ test space the transported
topology $(G_t)_*\eta$; the second map and the commutative square
are continuous by \eqref{eq:cq-testheat}. This specifies the exact
topologies rather than claiming an unproved topological embedding
of a quotient image equipped with its subspace topology.

There is a fully explicit inverse-domain calculation in the same
test coordinates. For $\lambda\in\cqC$, $k\geq1$, the image is
\begin{equation}\label{eq:cq-testmoments}
 (T_t-\lambda)^k\mathcal E
 =\left\{g\in\mathcal E:
 \int_{\mathbb R}v^j e^{-tv^2/4-i\lambda v}g(v)\,dv=0
                  \quad(0\leq j<k)\right\}.
\end{equation}
On that precise image the inverse is
\begin{equation}\label{eq:cq-testinverse}
 [(T_t-\lambda)^{-k}g](u)
 =\frac{i^ke^{tu^2/4+i\lambda u}}{(k-1)!}
 \int_{-\infty}^{u}(u-v)^{k-1}
            e^{-tv^2/4-i\lambda v}g(v)\,dv.
\end{equation}
To prove both directions, write
$f=e^{tu^2/4+i\lambda u}h$.
Direct differentiation gives
$(T_t-\lambda)f=e^{tu^2/4+i\lambda u}(-i h')$.
Iteration reduces the inverse equation to a $k$th primitive, giving
\eqref{eq:cq-testinverse} because $(-i)^ki^k=1$.
For $g$ supported in $[a,b]$, that formula vanishes below $a$.
Above $b$, expansion of $(u-v)^{k-1}$ shows it vanishes identically
exactly when all $k$ moments in \eqref{eq:cq-testmoments} vanish;
the binomial coefficients and the nonzero exponential do not change
that equivalence. Its smoothness follows by repeated differentiation
of the integral; smoothness of $g$ with compact support gives
matching derivatives at both endpoints. Necessity for any compact
solution also follows by integrating $h^{(k)}$ against
$1,v,\ldots,v^{k-1}$ and integrating by parts $k$ times.
Uniqueness follows since a compactly supported solution of
$h^{(k)}=0$ is a polynomial of degree at most $k-1$ which vanishes
on a ray, and hence is zero. This proves the exact image, inverse,
kernel zero, and preservation of the original support interval.

In particular the inverse of $\mathcal R_t^{\rm test}$ is
\[
 (g_0,g_1)\longmapsto(g_1,-\tfrac12T_t^{-1}g_0),
 \qquad
 \operatorname{Dom}(\mathcal R_t^{\rm test})^{-1}
 =\left\{(g_0,g_1):\int e^{-tv^2/4}g_0(v)\,dv=0\right\}.
\]
Substitution into its two rows proves both inverse identities.
Thus the branch-locus obstruction in the concrete test realization
is a retained Gaussian-weighted zero-moment condition. The exact
arithmetic quotient instead has invertible $S_t$ and root operator,
as proved above. Their ingress map and its kernels explain precisely
how these statements coexist; one operator's inverse domain is not
silently assigned to the other realization.

\section{Proof scope and reproducibility}
The new source-topology assertions proved here are the complete
transported closure, homeomorphic quotient, conjugated spectral
operator, source-domain differential correspondence, exact coordinate
kernels and cokernels, correlated jet extension, finite free cover,
compact-test heat ingress with its weighted-moment inverse domains,
and exact comparison diagrams. The compact and growth relation
subspaces both equal $\Xi\cqE$ by the complete division proof.
Their quotient topologies are computed: the compact quotient has
its exact jet-image topology and product completion, whereas the
growth quotient is complete and the comparison inverse is not
continuous. All derivative relation domains and finite ambiguity
projections retain the full zero multiplicities and analytic units.
The arithmetic principal-part realization has the exact kernel
$\Xi\cqE$, every finite image, the same product completion,
and explicit original-coordinate spectral blocks. The full Schwartz
trace factorization and its order-class intersection are proved
without assigning every Schwartz trace to that order class.
Both explicit realizations have proved frozen differential and
heat noninvariance and the original-coordinate divisor spectrum.
Each heated finite Hermite trace also has its proved unique
M\"obius inverse, precise Mellin domain and full nontrivial-zero
pole orders, with actual failures of boundedness at the origin
for suitable original seeds at every nonzero time.
The missing identification of the paper's unnamed topology is
recorded explicitly and is not replaced by a conditional arithmetic
theorem. Nothing in these results establishes an off-critical zero,
a negative Weil value, or an arithmetic endpoint by analogy with
the inverse-fibre collision. The exact transport through that
collision cover and the derivative information it requires have
been calculated above.

The companion checker uses rational polynomial arithmetic to verify
the heat conjugation on finite polynomial inputs, the retained
product terms, the matrix commutator, the iterated branch-domain
coefficients, the original-coordinate relation, and sample finite
jet multiplication, the multiplicity-sensitive derivative domains,
the source heat commutator recurrence, exact M\"obius divisor sums,
and Mellin principal-part coordinate inverses.
These are reproducible checks of the displayed
identities, not a substitute for their written proofs. Source
hashes and exact reading locators are in the companion provenance
receipt; no numerical zero search or Lean run is part of this work.
